\chapter{Експериментальний додаток}
\label{sec:supplement}

Для кожного розділу тут зібрано його головні твердження і те, на чому вони тримаються: дослід, у якому це виміряно, що саме виміряно і де це опубліковано. Статус у правому стовпці показує, наскільки твердження тверде: \emph{виміряно} --- є прямий дослід; \emph{теорія} --- математичний наслідок, виведений у розділі або в класичній роботі; \emph{модель} --- результат розрахунку за моделлю книги (скрипти в теці \texttt{models/}); \emph{оцінка} --- число порядку величини без прямого вимірювання; \emph{гіпотеза} --- правдоподібне пояснення, яке дослідом ще не доведено. Де досліду немає, так і сказано.

\section{Розділ~\ref{sec:neurontypes}. Типи нейронів}

Числа розділу спираються на прямі підрахунки клітин у мозку людини там, де такі підрахунки є, правило «один нейрон --- один медіатор» --- на клітинні атласи й досліди, які знайшли й винятки, роль \nt{DOPA} --- на записи імпульсів, а ролі решти широкомовних каналів --- на гіпотези.

{\footnotesize
\setlength{\tabcolsep}{3pt}\renewcommand{\arraystretch}{1.15}
\begin{longtable}{>{\raggedright}p{0.5cm}>{\raggedright}p{4.0cm}>{\raggedright}p{5.6cm}>{\raggedright}p{2.3cm}>{\raggedright\arraybackslash}p{1.7cm}}
\toprule
№ & Твердження & Дослід і що виміряно & Джерело & Статус \\
\midrule\endfirsthead
\toprule
№ & Твердження & Дослід і що виміряно & Джерело & Статус \\
\midrule\endhead
1 & У мозку людини близько $8.6\cdot10^{10}$ нейронів, і майже всі \nt{GLUT}-нейрони --- дрібні нейрони мозочка (табл.~\ref{tab:types}) & Ізотропний фракціонатор: тканину розчиняють в однорідну суспензію ядер і рахують ядра з нейронним маркером NeuN. У дорослих чоловіків $86.1\pm8.1$ млрд нейронів; у корі великих півкуль із них лише $19\,\%$, основна маса --- в мозочку & \cite{Azevedo2009} & виміряно \\
2 & Майже кожен нейрон має один головний медіатор (твердження~\ref{prop:dale}), але є винятки & Транскриптомний атлас мозку миші: близько $4\cdot10^6$ клітин, $5322$ кластери; кластерам приписано медіатор за генами синтезу й пакування. Винятки виміряно безпосередньо: \nt{DOPA}-нейрони у зрізах викидають ще й GABA через той самий везикулярний транспортер і гальмують нейрони стріатума; у частини \nt{DOPA}-аксонів глутамат і дофамін виходять із різних закінчень одного проводу (електронна мікроскопія та оптогенетика) & \cite{Strata1999,Yao2023,Tritsch2012,Zhang2015} & виміряно \\
3 & Увесь стовпець матриці ваг одного знака~\eqref{eq:dalesign} & Висновок у розділі з твердження~\ref{prop:dale} і того, що рецептори глутамату майже всі збуджувальні, а GABA --- гальмівні. Прямої перевірки «усі отримувачі одного нейрона отримують вагу одного знака» як загального правила немає; контрприклади --- спільний викид GABA \nt{DOPA}-нейронами (рядок~2) і отримувачі з рецепторами різного знака (рядок~11) & висновок у розділі & теорія \\
4 & Від трьох чвертей до чотирьох п'ятих нейронів кори --- \nt{GLUT}, від п'ятої частини до чверті --- \nt{GABA} & Імунофарбування на GABA у стовпчиках завширшки $50$\,мкм крізь усю товщу кори в $10$ областях кори мавпи: у $8$ областях GABA-нейрони становлять близько $25\,\%$ усіх нейронів, у первинній зоровій корі --- менше $20\,\%$ & \cite{Hendry1987} & виміряно \\
5 & \nt{GLUT}-нейрон має близько $10^4$ виходів & Стереологія на розтині: у новій корі молодих чоловіків $1.64\cdot10^{14}$ синапсів (електронна мікроскопія, дисектор) і близько $2.3\cdot10^{10}$ нейронів. Відношення дає $\sim7\cdot10^3$ синапсів на нейрон; у середньому кількість входів дорівнює кількості виходів & \cite{Tang2001synapses,Pakkenberg1997} & виміряно \\
6 & Вихід \nt{DOPA}-нейрона розгалужується на $10^5$--$10^6$ закінчень; це оцінка за охопленням мішені, а не підрахунок & Вірусна мітка мембрани окремих \nt{DOPA}-нейронів щура й повна реконструкція аксона: один нейрон покриває від $0.45$ до $5.7\,\%$ (у середньому $2.7\,\%$) об'єму стріатума. Виміряно охоплення, а не кількість закінчень: її виведено з охоплення за порядком величини, прямого підрахунку закінчень немає. Для \nt{SERO} і \nt{NORA} кількості закінчень --- грубі оцінки & \cite{Matsuda2009} & виміряно \\
7 & Широкомовних нейронів мало: \nt{DOPA} $\sim5\cdot10^5$ (головна з двох груп, нижня межа), \nt{SERO} $\sim3\cdot10^5$ (сума двох часткових підрахунків), \nt{HIST} $\sim6\cdot10^4$ (табл.~\ref{tab:types}, рис.~\ref{fig:typepie}) & Стереологічні підрахунки в людини: $550\,000$ пігментованих нейронів у чорній субстанції (без вентральної покришки); $235\,000$ нейронів у дорсальному ядрі шва (усі нейрони ядра) і ще $125\,000$ серотонінових нейронів у покришці мосту; близько $64\,000$ гістамінових нейронів гіпоталамуса. Для \nt{NORA}, \nt{GLYC}, \nt{ACET} і \nt{ADRE} прямого підрахунку немає: у таблиці це грубі оцінки за порядком величини & \cite{Pakkenberg1991,Baker1990,Baker1991,Airaksinen1991} & виміряно \\
8 & Глутамат у щілині злітає приблизно до мілімоля й спадає за $\sim1\ms$ & За витісненням швидко від'єднуваного конкурентного блокатора з рецепторів NMDA під час синаптичної передачі (культура нейронів гіпокампа): пік $1.1$\,мМ, стала спаду $1.2\ms$ & \cite{Clements1992} & виміряно \\
9 & У працюючій корі збуджувальний і гальмівний струми майже врівноважені, нейрон сидить біля порога & Теорія: мережа із сильними зв'язками сама приходить у врівноважений режим. Дослід: у префронтальній корі тхора in vivo під час «верхніх» станів потенціал реверсії синаптичного струму тримається біля $-37$\,мВ сотні мілісекунд, хоча провідність змінюється на $\sim21$\,нСм: збудження й гальмування зростають і спадають разом & \cite{vanVreeswijk1996,Haider2006} & виміряно \\
10 & \nt{DOPA}-нейрони повідомляють помилку передбачення винагороди~\eqref{eq:rpe} (рис.~\ref{fig:rpe}) & Імпульси \nt{DOPA}-нейронів мавпи: пачка на несподіваний сік; після навчання --- пачка на сигнал і немає відповіді на передбачений сік; провал частоти, коли обіцяного соку не дано. У кількісному досліді частота пропорційна різниці поточної винагороди та зваженого середнього минулих, але лише для додатних помилок. Саме рівняння~\eqref{eq:rpe} --- алгоритм часових різниць & \cite{Schultz1997,Bayer2005,SuttonBarto2018} & виміряно \\
11 & Знак і швидкість задає рецептор: іонотропні --- частки мілісекунди, метаботропні --- набагато повільніше (твердження~\ref{prop:receptor}) & Швидкі імпульси глутамату на клаптики мембрани нейронів гіпокампа: струм через іонотропні рецептори наростає (від $20$ до $80\,\%$) за $0.2$--$0.6\ms$ і спадає за $2$--$3\ms$. Повільний гальмівний потенціал пірамідних нейронів знімається вибірковим блокатором метаботропного рецептора GABA. Дві родини рецепторів дофаміну з протилежною дією; у стріатумі нейронам із рецептором D1 дофамін потрібен для посилення синапсів, нейронам із D2 --- для послаблення & \cite{Colquhoun1992,DutarNicoll1988,Beaulieu2011,Shen2008} & виміряно \\
12 & Широкомовний канал --- не один провід, а джгут із небагатьох ліній (підрозділ~\ref{sec:buses}) & Вірусно-генетичне маркування серотонінових нейронів дорсального ядра шва миші: нейрони, що посилають вихід у мигдалину й у лобову кору, лежать у різних місцях ядра, розгалужуються в області, які доповнюють одна одну, і протилежно відповідають на неприємний подразник. Огляд: підгрупи нейронів блакитної плями (\nt{NORA}) кодують різні процеси & \cite{Ren2018,Poe2020} & виміряно \\
13 & \nt{NORA} задає коефіцієнт підсилення $g$ & Гіпотеза адаптивного підсилення; модель мережі, у якій катехоламіни підвищують крутість передавальної характеристики. Прямого вимірювання «$g$ зростає з норадреналіном» у всій мережі немає; виміряно, що діаметр зіниці слідує за імпульсами нейронів блакитної плями мавпи & \cite{AstonJones2005,ServanSchreiber1990,Joshi2016} & гіпотеза \\
14 & \nt{ACET} перемикає «запис/читання» і задає швидкість навчання & Гіпотеза. Опора: у зрізах нюхової кори щура агоністи ацетилхоліну ($100$\,мкМ) послаблюють внутрішні, «пригадувальні» синапси більш ніж на $60\,\%$, а вхідні --- менш ніж на $15\,\%$ & \cite{Hasselmo2006,HasselmoBower1992} & гіпотеза \\
15 & \nt{SERO} задає коефіцієнт дисконтування $\gamma$; серотонінові нейрони працюють рівно, кілька імпульсів за секунду, і майже мовчать в одній із фаз сну & Гіпотеза «серотонін $=\gamma$». Найсильніший аргумент: оптогенетична активація серотонінових нейронів дорсального ядра шва миші зменшує кількість передчасних відмов від очікування відкладеної винагороди. Частоту імпульсів і мовчання уві сні зі швидкими рухами очей виміряно (огляд записів) & \cite{Doya2002,Miyazaki2014,JacobsAzmitia1992} & гіпотеза \\
\bottomrule
\end{longtable}}

\section{Розділ~\ref{sec:glutcomputer}. Що обчислюють \nt{GLUT}-нейрони}

Логічні твердження розділу --- теореми про схеми з невід'ємних ваг, виведені в самому розділі; дослід підтверджує модель нейрона й те, що схеми, зібрані за цими теоремами (детектор збігів, лінія затримки, самопідтримувана група, ланцюжок), у мозку справді трапляються.

{\footnotesize
\setlength{\tabcolsep}{3pt}\renewcommand{\arraystretch}{1.15}
\begin{longtable}{>{\raggedright}p{0.5cm}>{\raggedright}p{4.0cm}>{\raggedright}p{5.6cm}>{\raggedright}p{2.3cm}>{\raggedright\arraybackslash}p{1.7cm}}
\toprule
№ & Твердження & Дослід і що виміряно & Джерело & Статус \\
\midrule\endfirsthead
\toprule
№ & Твердження & Дослід і що виміряно & Джерело & Статус \\
\midrule\endhead
1 & Нейрон --- RC-коло з порогом і скиданням~\eqref{eq:glutneuron} & У тіло пірамідних нейронів кори щура (зрізи) подавали шумовий струм, схожий на синаптичний потік у живому мозку. Залежність середньої частоти від середнього й розкиду струму описується інтегрувальним нейроном із витоком, якщо додати повільну адаптацію частоти. Діапазони $\tau$ і затримок наведено в розділі без посилання на дослід & \cite{Rauch2003} & виміряно \\
2 & Модель~\eqref{eq:glutneuron} --- спрощення: гілки входів самі видають місцеві імпульси & Запис безпосередньо з дендритів пірамідних нейронів зорової кори миші in vivo: зоровий подразник викликає місцеві дендритні імпульси, що залежать від рецепторів NMDA; їх придушення розширює налаштування нейрона на орієнтацію & \cite{Smith2013} & виміряно \\
3 & Вікно збігу $\Delta_{\max}=\tau\ln\frac{w}{\theta-w}$~\eqref{eq:window}; при $\theta=1.5w$ воно дорівнює $0.7\tau$ & Наслідок моделі~\eqref{eq:glutneuron}. Пряме вимірювання вузького вікна підсумовування є для нейронів зі швидким гальмуванням (розділ~\ref{sec:gaba}, рядок~3 його таблиці) & висновок у розділі & теорія \\
4 & Мережа з самих лише \nt{GLUT}-нейронів обчислює лише монотонні функції входів (твердження~\ref{prop:monotone}) & Доведення в розділі: ітерація від тиші з неспадною правою частиною. Це твердження про модель; досліду, який перевіряв би його на живій мережі без гальмування, немає & висновок у розділі & теорія \\
5 & Органи чуття дають пару проводів: одні клітини відповідають на посилення подразника, інші --- на послаблення & Сітківка: уже на біполярних клітинах сигнал колбочок розходиться на вмикальний і вимикальний канали. Фармакологічний блок вмикальних біполярних клітин у мавпи: канали йдуть окремо до кори; після блоку погіршується виявлення посилення світла, але не послаблення & \cite{Schiller1992} & виміряно \\
6 & З двопровідним кодом мережа з \nt{GLUT}-нейронів обчислює будь-яку логічну функцію асинхронно, без спільного такту, ціною вдвічі більшої кількості нейронів (твердження~\ref{prop:dualrail}, \eqref{eq:dualrail}, рис.~\ref{fig:dualxor}) & Двопровідний код і визначення готовності за самим сигналом --- стандартна техніка асинхронних схем, нечутливих до затримок. Схему виключного АБО з шести нейронів зібрано в розділі. Досліду, який показував би, що мозок обчислює логічні функції саме двопровідним кодом, немає & \cite{Sparso2001} & теорія \\
7 & Найбільша різниця часу приходу звуку у вуха в людини $\approx0.6\ms$, а мозок розрізняє близько десяти мікросекунд & Слухачі в досліді з двома варіантами відповіді: поріг розрізнення різниці часу (75\,\% правильних) --- $9$\,мкс для шуму в смузі $150$--$1700$\,Гц, $11$\,мкс для тону $1$\,кГц, $28$\,мкс для клацання. Межа $0.6\ms$ --- розрахунок розділу, $0.22\,\text{м}/343\,\text{м/с}$ & \cite{KlumppEady1956} & виміряно \\
8 & У птахів різниця часів перетворюється на місце через лінії затримки й детектори збігів~\eqref{eq:itd} (рис.~\ref{fig:itd}) & Сипуха: аксони від двох вух входять у ядро детекторів збігів із протилежних боків; час приходу імпульсів уздовж аксона змінюється впорядковано з глибиною, і розкид затримок крізь ядро ($160$\,мкс на $700$\,мкм) збігається з діапазоном міжвушних затримок & \cite{Carr1990} & виміряно \\
9 & У ссавців ряду затримок не знайшли; ту саму задачу розв'язують детектори збігів із точно розрахованим гальмуванням від \nt{GLYC}-нейронів & Запис in vivo в медіальній верхній оливі піщанки: відповіді не утворюють карти напрямків; підведення гліцину та його блокатора через мікропіпетку показує, що точно розраховане в часі \emph{гліцинове} гальмування задає робочий діапазон затримок. Гальмування тут від \nt{GLYC}, а не від \nt{GABA} & \cite{Brand2002,Grothe2010} & виміряно \\
10 & Група із сильним зворотним зв'язком --- тригер; самопідтримувана активність триває секунди після подразника (рис.~\ref{fig:bistable}) & Префронтальна кора мавпи в задачі з відкладеним рухом очей до запам'ятованої точки (затримка $1$--$6$\,с): $87$ із $170$ нейронів, пов'язаних із задачею, змінюють частоту під час затримки, у $79\,\%$ із них це залежить від напрямку на запам'ятовану точку. Що цю активність тримає зворотний зв'язок із двома стійкими станами, --- теорія & \cite{Funahashi1989,Wang2001} & виміряно \\
11 & Біт записується, якщо приплив триває довше $\approx0.7\tau$ (рис.~\ref{fig:recurrentgroup}б) & Розрахунок рівняння $\tau\,\dd r/\dd t=-r+f(wr+I)$ методом Ейлера при $w=1$, $I=0.6$; скрипту в \texttt{models/} немає. Досліду немає & розрахунок у розділі & модель \\
12 & Пам'ять можна зберігати в полегшенні синапсів, майже без імпульсів & Теорія: залишковий кальцій у закінченнях тримає запис близько секунди. Опора: у медіальній префронтальній корі тхора (запис із кількох нейронів одночасно) є підмережа пірамідних нейронів, зв'язаних полегшуваними синапсами з тривалою післядією. Огляд спостережень «тихих» проміжків під час утримання в пам'яті. Який спосіб кора використовує коли --- відкрите питання & \cite{Mongillo2008,WangMarkram2006,Stokes2015} & гіпотеза \\
13 & Пам'ять стирається втомою: запас медіатора виснажується & Парні записи пірамідних нейронів кори: відповідь синапсу падає за частих імпульсів, швидкість падіння задається ймовірністю викиду. Що саме це гасить самопідтримувану групу, безпосередньо не показано & \cite{Tsodyks1997} & виміряно \\
14 & Ланцюжок груп --- таймер, що запускається подією; у співочих птахів нейрони спрацьовують суворо по черзі & Запис упізнаних нейронів вищого вокального центру зебрової амадини уві сні й під час співу: кожен нейрон, що посилає вихід у моторне ядро, видає одну коротку пачку в один точний момент пісні, і нейрони спрацьовують послідовно & \cite{Hahnloser2002} & виміряно \\
\bottomrule
\end{longtable}}

\section{Розділ~\ref{sec:gaba}. \nt{GLUT} і \nt{GABA} разом}

Логічні й динамічні твердження розділу виведено в ньому самому з лінійного аналізу й закону Ома; дослід показує, що схеми, на які вони спираються, --- заборона із затримкою, пряме гальмування слідом за збудженням, стабілізація гальмуванням, ритм петлі \nt{GLUT}--\nt{GABA}, --- у мозку справді працюють.

{\footnotesize
\setlength{\tabcolsep}{3pt}\renewcommand{\arraystretch}{1.15}
\begin{longtable}{>{\raggedright}p{0.5cm}>{\raggedright}p{4.0cm}>{\raggedright}p{5.6cm}>{\raggedright}p{2.3cm}>{\raggedright\arraybackslash}p{1.7cm}}
\toprule
№ & Твердження & Дослід і що виміряно & Джерело & Статус \\
\midrule\endfirsthead
\toprule
№ & Твердження & Дослід і що виміряно & Джерело & Статус \\
\midrule\endhead
1 & \nt{GABA}-нейронів у корі від п'ятої частини до чверті & Імунофарбування на GABA у $10$ областях кори мавпи: близько $25\,\%$ нейронів у $8$ областях, менше $20\,\%$ у первинній зоровій корі. Огляд різноманіття гальмівних нейронів кори & \cite{Hendry1987,Markram2004} & виміряно \\
2 & Що робить гальмування, багато в чому вирішує місце посадки: дендрит, тіло, місце народження імпульсу, закінчення чужого проводу & Огляд: класи \nt{GABA}-нейронів кори різняться мішенню на отримувачі. Одночасний запис із тіла й дендрита пірамідного нейрона: пряме гальмування набагато сильніше на тілі, ніж на дендритах. Гальмування на закінченні чужого проводу в спинному мозку зменшує викид медіатора однієї вхідної лінії (огляд вимірювань) & \cite{Markram2004,Pouille2001,Rudomin1999} & виміряно \\
3 & Зміна знака через посередника коштує однієї затримки, близько мілісекунди; гальмування, що прийшло слідом за збудженням, звужує вікно збігу & Пірамідні нейрони гіпокампа щура у зрізах: гальмування через одного посередника приходить слідом за прямим збудженням із короткою затримкою, і вікно, у якому два збуджувальні входи підсумовуються до порога, менше за $2\ms$. На дендритах, де це гальмування слабше, вікно ширше & \cite{Pouille2001} & виміряно \\
4 & Заборона $a\wedge\bar b$~\eqref{eq:veto} з АБО й сталою одиницею функціонально повна (твердження~\ref{prop:gabacomplete}) & Доведення в розділі. Класичний результат: мережа порогових елементів із гальмівними входами реалізує будь-який логічний вираз. Твердження про модель, досліду немає & \cite{McCullochPitts1943} & теорія \\
5 & Заборона із затримкою дає детектор напрямку руху з оптимальною швидкістю $v^*=\Delta x/d$~\eqref{eq:vstar} & Сітківка кроля: вибірковість до напрямку створюється гальмуванням, яке за руху в «нульовому» напрямку гасить збудження. Окремий запис збуджувального й гальмівного входів вибіркових клітин показав, що зірчасті амакринові клітини (\nt{GABA}) гальмують їх несиметрично --- лише відростками, спрямованими в «нульовий» бік. Формула $v^*$ --- висновок розділу & \cite{Barlow1965,Fried2002} & виміряно \\
6 & Шунтувальне гальмування ділить напругу, а не віднімає~\eqref{eq:shunt} (рис.~\ref{fig:shunt}) & Закон Ома для подільника з трьох провідностей за рівноважного потенціалу хлору біля спокою; для нейрона без імпульсів & висновок у розділі & теорія \\
7 & На частоту імпульсів шунт діє відніманням, а ділення виходить від гальмування разом зі врівноваженим фоновим шумом & Модель: у нейрона, що спрацьовує, струм через шунт майже сталий, і ефект на частоту віднімальний. Дослід: у пірамідні нейрони соматосенсорної кори щура (зрізи) через динамічну фіксацію вводили фонові збуджувальні й гальмівні провідності; одночасне врівноважене зростання обох ділить крутість відповіді без помітного зсуву частоти. Огляд: ділення на загальну активність трапляється в багатьох відділах мозку & \cite{HoltKoch1997,Chance2002,Carandini2012} & виміряно \\
8 & При $\beta>1$ взаємне гальмування обирає одного переможця (твердження~\ref{prop:wta}, \eqref{eq:wta}) & Власні значення $-(1\pm\beta)/\tau$ --- висновок у розділі. Частоти $0.73$ і $0.53$ при $\beta=0.5$ --- нерухома точка $r_{1,2}=(I_{1,2}-\beta I_{2,1})/(1-\beta^2)$; скрипту в \texttt{models/} немає. Прямого досліду в розділі не наведено & висновок у розділі & теорія \\
9 & Скидання пам'яті сигналом через \nt{GABA}; дві групи із взаємним гальмуванням --- засувка & Випливає з рис.~\ref{fig:bistable} і твердження~\ref{prop:wta}. Досліду немає & висновок у розділі & теорія \\
10 & Рівновага петлі \nt{GLUT}--\nt{GABA} стійка, якщо гальмування досить сильне й досить швидке (твердження~\ref{prop:ei}) & Модель двох популяцій~\eqref{eq:ei}; умови (i) і (ii) --- визначник і слід її матриці, виведені в розділі & \cite{WilsonCowan1972} & теорія \\
11 & При $w_{EE}=1.5$ і $w_{EI}w_{IE}=2$ швидке гальмування ($\tau_I=2\ms$) тримає $E^*=0.67h$, повільне ($30\ms$) розгойдує коливання (рис.~\ref{fig:ei}) & Чисельний розв'язок~\eqref{eq:ei}, скрипт \texttt{figures/make\_ei.py} & розрахунок & модель \\
12 & Кора працює в режимі стабілізації гальмуванням; підпис --- підштовхнеш \nt{GABA}-нейрони, і їхня частота впаде~\eqref{eq:isn} & Теорія передбачила парадоксальну відповідь. Дослід: внутрішньоклітинні записи в первинній зоровій корі кішки --- при пригніченні відповіді подразником навколо гальмівна провідність не зростає, а падає разом зі збуджувальною. Оптогенетичне збудження PV-нейронів миші в зоровій, соматосенсорній і моторній корі знижує частоту гальмівних нейронів, зокрема без подразника й за легкого наркозу. Коефіцієнт $-1/3$ за чисел рис.~\ref{fig:ei} --- висновок розділу & \cite{Tsodyks1997isn,Ozeki2009,Sanzeni2020} & виміряно \\
13 & Петля «\nt{GLUT} збуджує \nt{GABA}, \nt{GABA} гальмує \nt{GLUT}» породжує швидкий ритм із періодом $12$--$50\ms$ & Оптогенетична стимуляція швидких \nt{GABA}-нейронів у соматосенсорній корі миші на частотах $8$--$200$\,Гц вибірково посилює гамма-ритм ($20$--$80$\,Гц, період $12$--$50\ms$); стимуляція \nt{GLUT}-нейронів посилює лише повільніші ритми & \cite{Cardin2009,Buzsaki2012} & виміряно \\
\bottomrule
\end{longtable}}

\section{Розділ~\ref{sec:code}. Абетка Морзе}

Граничні оцінки розділу --- теорія інформації та рахунок; виміряно, де в каналі виникає шум, як швидко працює мозок і скільки коштує імпульс, а питання про те, яким кодом кора користується насправді, лишається відкритим.

{\footnotesize
\setlength{\tabcolsep}{3pt}\renewcommand{\arraystretch}{1.15}
\begin{longtable}{>{\raggedright}p{0.5cm}>{\raggedright}p{4.0cm}>{\raggedright}p{5.6cm}>{\raggedright}p{2.3cm}>{\raggedright\arraybackslash}p{1.7cm}}
\toprule
№ & Твердження & Дослід і що виміряно & Джерело & Статус \\
\midrule\endfirsthead
\toprule
№ & Твердження & Дослід і що виміряно & Джерело & Статус \\
\midrule\endhead
1 & Частоти \nt{GLUT}-нейронів кори --- від часток до десятків імпульсів за секунду, і більшу частину часу нейрони працюють біля нижньої межі & Запис через піпетку, прикладену до клітини (вибір нейрона не залежить від його активності), у слуховій корі щура, що не спить: у кожен момент частоту вище $20$ імпульсів за секунду дають менше $5\,\%$ нейронів; у частини нейронів фонова частота нижча за $0.01$ за секунду; розподіл частот логнормальний & \cite{Hromadka2008} & виміряно \\
2 & Пропускна здатність імпульсного каналу $R_{\max}\approx r\log_2\frac{e}{r\Delta t}$~\eqref{eq:spikeentropy}, майже вся вона --- у часі імпульсів (твердження~\ref{prop:capacity}) & Ентропія двійкового рядка з клітинок довжини $\Delta t$. Числа розділу: $8$ біт на імпульс при $r=10$ за секунду й $\Delta t=1\ms$, близько $5$ біт при $\Delta t=10\ms$, менше $2$ біт за секунду в лічильника імпульсів & \cite{MacKay1952} & теорія \\
3 & Чутливі нейрони передають від одного до кількох біт на імпульс, у найкращих випадках до половини межі & Відновлення подразника за імпульсами чутливих нейронів, для яких подразник відомий; зведення вимірювань у книзі & \cite{Rieke1997} & виміряно \\
4 & Генератор імпульсів точний; точний час задають швидкі перепади входу & Нейрони кори щура у зрізах: на повторюваний струм зі швидкими коливаннями імпульси повторюються з точністю менше $1\ms$, на сталий струм моменти розповзаються & \cite{Mainen1995} & виміряно \\
5 & Синапс ненадійний: понад половина імпульсів не доходить до отримувача через випадковість викиду & Мінімальна стимуляція входів пірамідних нейронів гіпокампа (зрізи): понад половина імпульсів не викликає відповіді. Виключено коливання порога стимуляції, збої проведення на розгалуженнях аксона й низьку температуру досліду; лишається ймовірнісний викид & \cite{AllenStevens1994} & виміряно \\
6 & Потік імпульсів у живій корі виглядає випадковим: інтервали розкидані як у пуассонівського потоку, дисперсія кількості імпульсів близька до середнього; частина «шуму» --- повідомлення про рухи тварини & Записи в зорових областях V1 і MT мавпи, що не спить: коефіцієнт варіації інтервалів близько $1$, відношення дисперсії кількості імпульсів до середнього --- як у випадкового процесу; модель, що підсумовує незалежні малі входи, дає розкид набагато менший. У зоровій корі миші значну частину розкиду передбачають за відеозаписом рухів самої тварини & \cite{Softky1993,Stringer2019} & виміряно \\
7 & Частотний код повільний: похибка $1/\sqrt{rT}$~\eqref{eq:rateerror}, а мозок упізнає картинку за $\sim150\ms$ & Формула --- пуассонівська статистика, висновок розділу. Дослід: люди вирішували, чи є тварина на незнайомій фотографії, показаній на $20\ms$; потенціали мозку на «так» і «ні» розходяться приблизно через $150\ms$. Розбиття цього часу на десяток ступенів по $10$--$20\ms$ --- оцінка розділу, не вимірювання & \cite{Thorpe1996} & виміряно \\
8 & Усереднення за групою обмежене кореляцією: похибка падає не більше ніж у $1/\sqrt c$ разів~\eqref{eq:correlated} (твердження~\ref{prop:pooling}) & Пари нейронів області MT мавпи, записані одночасно: коефіцієнт кореляції кількостей імпульсів у середньому $0.12$, і теоретичний аналіз показує, що навіть така кореляція обмежує виграш від групи. Теорія: межу ставить лише та частина кореляції, що схожа на сигнал. Модель протилежної позиції: частоту групи з $50$--$100$ нейронів читають за один міжімпульсний інтервал, $10$--$50\ms$, і час окремих імпульсів у корі не використовується & \cite{Zohary1994,MorenoBote2014,ShadlenNewsome1998} & виміряно \\
9 & Число можна писати в час першого імпульсу~\eqref{eq:latency}, у порядок імпульсів групи або у фазу місцевого ритму & Сітківка саламандри: просторову структуру коротко показаної картинки записано у відносних затримках перших імпульсів вихідних нейронів, майже незалежно від контрасту. Клітини місця гіпокампа щура: під час проходу через «своє» місце імпульси зсуваються дедалі раніше в циклі тета-ритму ($7$--$12$\,Гц), зсув від $100^\circ$ до $355^\circ$. Наскільки кора користується часом імпульсів --- відкрите питання (рядок~8) & \cite{Gollisch2008,OKeefe1993} & виміряно \\
10 & Який код прочитано, вирішує стала часу отримувача~\eqref{eq:reader} (твердження~\ref{prop:reader}); в активній корі нейрони ближчі до детекторів збігів & Формула~\eqref{eq:reader} --- висновок розділу. Огляд записів in vivo: в активній корі провідність мембрани в кілька разів вища, ніж у спокої, і стала часу відповідно коротша. Що кора перемикає код саме так, безпосередньо не показано & \cite{Destexhe2003} & гіпотеза \\
11 & Синапс, що виснажується, передає зміну частоти, по суті $\Delta r/r$~\eqref{eq:depression} (рис.~\ref{fig:depression}) & Парні записи пірамідних нейронів кори: швидкість виснаження задається ймовірністю викиду; за повільного виснаження відповідь відображає частоту, за швидкого --- збіги. Модель на основі виміряного виснаження в корі щура: однакові відносні зміни частоти на швидких і повільних входах дають однакову відповідь, тобто синапс передає $\Delta r/r$. Числа $1.7\to2.2$, $6.7$ і $90\ms$ --- розрахунок розділу & \cite{Tsodyks1997,Abbott1997depression} & виміряно \\
12 & Пачка --- «тире»: її ймовірність загубитися $(1-p)^k$~\eqref{eq:burst} мала, полегшення зменшує її ще & Огляд: на ненадійних синапсах одиночні імпульси часто губляться, а пачки передаються надійно завдяки полегшенню; пачка викликає довготривалі зміни синапсів. Що мозок читає пачку як окремий знак, --- гіпотеза & \cite{Lisman1997,AllenStevens1994} & гіпотеза \\
13 & Швидкі \nt{GABA}-нейрони задають ритм і «пунктуацію»; ще один клас гальмує гальмівних і відчиняє ворота (твердження~\ref{prop:punctuation}) & Огляд властивостей класів \nt{GABA}-нейронів кори. Оптогенетика: ритмічне збудження швидких \nt{GABA}-нейронів викликає гамма-ритм, і відповідь на подразник залежить від фази циклу. У зоровій корі миші PV-нейрони гальмують один одного, SST-нейрони --- усі інші класи, VIP-нейрони --- переважно SST. Збудження VIP-нейронів у миші, що не спить, знімає гальмування з \nt{GLUT}-нейронів; їх вмикає сигнал підкріплення. Поділ ролей як загальне правило --- гіпотеза & \cite{Tremblay2016,Cardin2009,Pfeffer2013,Pi2013} & гіпотеза \\
14 & Енергія обмежує середню частоту величиною порядку одного імпульсу за секунду~\eqref{eq:energy} & Енергетичний бюджет сірої речовини гризуна за анатомічними й фізіологічними даними: імпульси й синаптичні струми --- $47\,\%$ і $34\,\%$ енергії сигналізації; одночасно можуть бути активними не більше $\sim15\,\%$ нейронів. Для кори людини: не більше $\sim1\,\%$ нейронів можуть бути сильно активними одночасно. Число $\sim10^9$ АТФ на імпульс і потужність $\sim10$\,Вт на сигналізацію --- оцінка розділу на основі цих розрахунків & \cite{Attwell2001,Lennie2003} & теорія \\
15 & Код розріджений і підлаштований під найбільшу кількість біт на джоуль; кора міняє точність на енергію (твердження~\ref{prop:economy}) & Гіпотеза ощадливого коду. Опора: у мишей за обмеження їжі в зоровій корі падає провідність AMPA-рецепторів і синаптична витрата АТФ ($-29\,\%$), частота імпульсів зберігається, а налаштування на орієнтацію розширюється на $32\,\%$ & \cite{Barlow1961,Padamsey2022} & гіпотеза \\
\bottomrule
\end{longtable}}

\section{Розділ~\ref{sec:serocomputer}. \nt{SERO}-нейрони}

Властивості самих \nt{SERO}-нейронів (ритм, автогальмування, знак за рецептором, підсистеми) виміряно прямо; обчислення розділу --- регулятор, фільтр, логіка --- випливають із його рівнянь; $\tau_c$, коефіцієнт дифузії серотоніну і кількість місць викиду --- оцінки; роль петлі як регулятора рівня в мозку --- гіпотеза (у ядрі шва автогальмування точкове і дає лише короткі паузи), а роль \nt{SERO} як горизонту --- гіпотеза, на користь якої свідчать поведінкові досліди.

{\footnotesize
\setlength{\tabcolsep}{3pt}\renewcommand{\arraystretch}{1.15}
\begin{longtable}{>{\raggedright}p{0.5cm}>{\raggedright}p{4.0cm}>{\raggedright}p{5.6cm}>{\raggedright}p{2.3cm}>{\raggedright\arraybackslash}p{1.7cm}}
\toprule
№ & Твердження & Дослід і що виміряно & Джерело & Статус \\
\midrule\endfirsthead
\toprule
№ & Твердження & Дослід і що виміряно & Джерело & Статус \\
\midrule\endhead
1 & Знак дії серотоніну вибирає рецептор отримувача (твердження~\ref{prop:receptor}) & Рецептори серотоніну поділено на родини за фармакологією і механізмом: 5-HT$_{1A}$ через G-білок відкриває калієві канали і гальмує клітину, 5-HT$_{2A}$ та іонотропний 5-HT$_3$ збуджують. Один медіатор дає протилежні відповіді в різних клітинах. & \cite{BarnesSharp1999} & виміряно \\
2 & \nt{SERO}-нейрон у стані неспання рівно видає кілька імпульсів на секунду & Запис нейронів дорсального ядра шва у вільно рухливих котів: повільний ритмічний розряд $2.8\pm0.2$ імп/с у спокійному неспанні; у повільному сні частота падає на $34$--$68\,\%$, у швидкому --- на $98\,\%$. Під наркозом у щурів --- $1.7\pm0.5$ імп/с, дуже регулярно. & \cite{TrulsonJacobs1979, GagnonParent2014, JacobsAzmitia1992} & виміряно \\
3 & \nt{SERO}-нейрон --- ритмоводій: поштовх на кожен імпульс не потрібен, але потрібен рівний фоновий приток (у зрізі --- норадреналін через $\alpha_1$-рецептори, в організмі --- від \nt{NORA}) & У зрізі мозку клітини розряджаються повільно і рівно, після кожного імпульсу --- слідова гіперполяризація на $200$--$800\ms$. Спонтанно активних клітин у зрізі менше, ніж в організмі: рівному ритму потрібен тонічний вхід норадреналіну через $\alpha_1$-рецептори, блокада яких його гасить. & \cite{VandermaelenAghajanian1983} & виміряно \\
4 & Майже всі рецептори метаботропні; відповідь повільна: наростає і спадає в межах приблизно секунди & Усі родини рецепторів, крім 5-HT$_3$, пов'язані з G-білками. У зрізі викликаний викид серотоніну дає через 5-HT$_{1A}$ гальмівний струм, який наростає і спадає в межах секунди. & \cite{BarnesSharp1999, CourtneyFord2016} & виміряно \\
5 & В областях призначення серотонін виходить в об'єм, а не лише в щілину; у самому ядрі шва гальмування сусідів іде точка-точка & Швидка циклічна вольтамперометрія у зрізах: викид на один імпульс однаковий в області із синапсами і в ядрі шва, де синапсів майже немає; не залежить від блокади захоплення і рецепторів; молекул на закінчення набагато менше, ніж переносників і рецепторів, а позаклітинна концентрація збігається зі спорідненістю головного рецептора. У ядрі шва гальмування через 5-HT$_{1A}$ іде точка-точка: переносник не дає серотоніну накопичитися поза щілиною. & \cite{Bunin1998, CourtneyFord2016} & виміряно \\
6 & Концентрація за~\eqref{eq:seroneuron} спадає через відкачування; $\tau_c$ порядку секунди --- оцінка & Вольтамперометрія в живому мозку: позаклітинний серотонін обмежують насамперед зворотне захоплення і руйнування, а не синтез. Прямого вимірювання $\tau_c$ у цитованих роботах немає; порядок величини --- оцінка розділу. & \cite{Hashemi2012serotonin} & виміряно; $\tau_c$ --- оцінка \\
7 & НЕ і АБО-НЕ на одному ритмоводії; повнота \nt{SERO}-логіки (твердження~\ref{prop:serocomplete}, рис.~\ref{fig:serogates}) & Випливає з~\eqref{eq:seroneuron}: гальмований ритмоводій --- АБО-НЕ, а АБО-НЕ функціонально повний. Досліду, де з \nt{SERO}-нейронів зібрано логіку, немає; умова окремих об'ємів у мозку не виконується. & виведено в розділі & теорія \\
8 & Серотонін гальмує сам \nt{SERO}-нейрон через рецептори на ньому & У щурів під наркозом агоністи 5-HT$_{1A}$ внутрішньовенно та іонофоретично гальмують розряд нейронів дорсального шва з тією самою залежністю доза--відповідь, що й сам серотонін; агоністи 5-HT$_{1B}$ майже не діють. У зрізі ендогенний серотонін сусідніх клітин дає через 5-HT$_{1A}$ струм і коротку паузу в розряді, не змінюючи середньої частоти. & \cite{Sprouse1987, CourtneyFord2016} & виміряно \\
9 & У моделі петля через спільний об'єм --- регулятор рівня: розкид і стала часу зменшуються в $1+G$ разів~\eqref{eq:feedback}; що петля так працює в мозку, --- гіпотеза & Класична формула підсилювача з від'ємним зворотним зв'язком; у розділі виведена заново. Підсилення петлі $G$ у \nt{SERO}-системи не виміряно. Виміряно: у зрізі автогальмування іде точка-точка, дає коротку паузу і не змінює середньої частоти. & \cite{Black1934feedback, CourtneyFord2016} & теорія; у мозку --- гіпотеза \\
10 & Концентрація --- ковзне середнє частоти, фільтр нижніх частот~\eqref{eq:lowpass}, рис.~\ref{fig:lowpass} & Розв'язок лінійного рівняння~\eqref{eq:seroneuron}; рисунок --- розрахунок за цією формулою при $\tau_c=1\,\text{с}$. Окремої моделі в \texttt{models/} немає. & виведено в розділі & теорія \\
11 & Звивистість щілин сповільнює дифузію малої молекули приблизно в $2.6$ раза & Метод точкового джерела: іонофорез тетраметиламонію і запис його концентрації іоноселективним електродом у зрізах і в живому мозку. Звивистість $\lambda\approx1.6$, тобто ефективний $D$ менший за вільний у $\lambda^2\approx2.6$ раза; частка міжклітинного об'єму близько $0.2$. Коефіцієнт дифузії самого серотоніну в тканині в цих роботах не вимірювали; у розділі він --- оцінка. & \cite{Sykova2008} & виміряно \\
12 & Довжина незалежного «проводу» $\ell\sim\sqrt{D\tau}\approx20$~мкм~\eqref{eq:diffusionlength}; кількість проводів обмежена кількістю таких областей (твердження~\ref{prop:volumewires}) & Закон дифузії і підстановка оцінок $D$ і $\tau$ (рядки~6 і~11); твердження~\ref{prop:volumewires} --- наслідок. Прямого досліду, який рахує незалежні канали об'ємної передачі, немає. & виведено в розділі & теорія \\
13 & Вихід одного \nt{SERO}-нейрона розкинутий по величезних ділянках мозку; місць викиду, за оцінкою, $\sim10^5$ & Повна реконструкція окремих нейронів дорсального шва в щура: усі висхідні аксони сильно розгалужуються, загальна довжина аксона однієї клітини до $18.7$~см, гілки однієї клітини сягають стріатума, префронтальної кори і мигдалеподібного тіла. Кількість місць викиду на клітину прямо не підраховано. & \cite{GagnonParent2014} & виміряно; $10^5$ --- оцінка \\
14 & \nt{SERO}-нейрони поділяються на кілька підсистем із різними виходами і відповідями & Вірусно-генетичне мічення в мишей: нейрони, що йдуть до мигдалеподібного тіла і до лобової кори, майже не перетинаються за розгалуженням, отримують різні входи і протилежно відповідають на неприємний стимул; увімкнення і вимкнення кожної групи по-різному змінює поведінку. & \cite{Ren2018} & виміряно \\
15 & Умова терпіння $D<\tau_\gamma\ln(R/r_0)$~\eqref{eq:patience} за експоненційного знецінення; за виміряного, гіперболічного, --- $D<\tau_\gamma(R/r_0-1)$ & Випливає зі знецінення $e^{-D/\tau_\gamma}$ або $1/(1+D/\tau_\gamma)$; виведено в розділі. Мавпи, вибір на затримках у секунди: знецінення ближче до гіперболи, ніж до експоненти; відповідь \nt{DOPA}-нейронів на провісник відкладеної винагороди спадає із затримкою теж за гіперболою. & виведено в розділі; \cite{KobayashiSchultz2008} & теорія \\
16 & \nt{SERO} задає горизонт $\tau_\gamma$ (підрозділ~\ref{sec:horizon}) & Оптогенетичне ввімкнення \nt{SERO}-нейронів дорсального шва в мишей подовжує очікування відкладеної винагороди і підвищує наполегливість у пошуку винагороди, що стала рідкісною. У людей зниження серотоніну (дієта без триптофану) робить знецінення відкладеної винагороди крутішим. Як концентрація перетворюється на $\tau_\gamma$, невідомо. & \cite{Doya2002, Miyazaki2014, Lottem2018, Schweighofer2008} & гіпотеза \\
\bottomrule
\end{longtable}}

\section{Розділ~\ref{sec:dofa}. Навчання з \nt{DOPA}}

Механізми синапсу (порції, детектор збігів, кальцій, вікна пластичності) і поведінку дофаміну як помилки передбачення виміряно прямо; те, що правила ведуть до градієнта і чого коштує широкомовлення, --- теореми; підсумок навчання вибору --- розрахунок за моделлю книги; схема, що обчислює помилку, --- гіпотеза.

{\footnotesize
\setlength{\tabcolsep}{3pt}\renewcommand{\arraystretch}{1.15}
\begin{longtable}{>{\raggedright}p{0.5cm}>{\raggedright}p{4.0cm}>{\raggedright}p{5.6cm}>{\raggedright}p{2.3cm}>{\raggedright\arraybackslash}p{1.7cm}}
\toprule
№ & Твердження & Дослід і що виміряно & Джерело & Статус \\
\midrule\endfirsthead
\toprule
№ & Твердження & Дослід і що виміряно & Джерело & Статус \\
\midrule\endhead
1 & Зворотного поширення помилки в тій формі, що в штучних мережах, у мозку не знайдено & Прямого досліду немає: це твердження про відсутність. Огляди розбирають, чого вимагає алгоритм (зворотні зв'язки, що повторюють прямі, окрема фаза) і які біологічні наближення до нього запропоновано. & \cite{Lillicrap2020, WhittingtonBogacz2019} & гіпотеза \\
2 & Вага розкладається на кількість місць викиду, імовірність викиду і відповідь на порцію: $w=N_s\,p\,A_1$~\eqref{eq:quantal} & Нервово-м'язовий синапс жаби за зниженого викиду: відповідь отримувача коливається сходинками, кратними спонтанній мініатюрній відповіді на одну порцію; кількість порцій на імпульс підлягає закону Пуассона. & \cite{delCastillo1954} & виміряно \\
3 & Рецептор отримувача --- детектор збігів; кальцій запускає зміну ваги & Патч-клемп нейронів миші: іони Mg$^{2+}$ закупорюють канал, відкритий глутаматом, за потенціалу спокою і виходять за деполяризації. Зрізи гіпокампа: хелатор кальцію в отримувачі блокує довготривале посилення, а вивільнення кальцію світлом усередині клітини саме посилює передачу. & \cite{Nowak1984, Malenka1988calcium} & виміряно \\
4 & Рішення виконує будь-який бік: росте $A_1$ в отримувача, змінюється $p$ у відправника & Глутамат, вивільнений світлом над одним шипиком, спричиняє швидкий і стійкий ріст шипика і струму через його рецептори; посилення супроводжується вбудовуванням AMPA-рецепторів. Деполяризація отримувача на певний час пригнічує викид у відправника; ефект переносять ендоканабіноїди, що йдуть назад через щілину (показано на \nt{GABA}-входах). Короткочасне виснаження залежить від імовірності викиду у відправника. & \cite{Matsuzaki2004, Malinow2002, WilsonNicoll2001, Tsodyks1997} & виміряно \\
5 & Синапс посилюється, коли відправник і отримувач активні разом~\eqref{eq:hebb} & Кролик під наркозом: після коротких серій імпульсів вхідним шляхом гіпокампа відповідь посилена на час від $30$~хв до $10$~год. У зрізі гіпокампа імпульси отримувача посилюють лише ті синапси, які в ту саму мить активні. & \cite{Bliss1973LTP, Kelso1986hebb} & виміряно \\
6 & Правило~\eqref{eq:oja} знаходить головний власний вектор кореляцій входів (твердження~\ref{prop:oja}) & Теорема; доведення в розділі. Досліду, де нейрон вивчив головну компоненту своїх входів, немає; механізм обмеження росту ваг у синапсі не встановлено. & \cite{Oja1982} & теорія \\
7 & Гебб за часом: вікно близько $\pm20\ms$ (рис.~\ref{fig:stdp}); з повторюваних послідовностей виростають ланцюжки & Пари нейронів гіпокампа в культурі: імпульс отримувача в межах $20\ms$ після імпульсу відправника дає стійке посилення, у межах $20\ms$ до --- послаблення; обидва залежать від NMDA-рецепторів. Відношення $A_-=0.6A_+$ на рисунку --- ілюстрація. Що так у мережі виростають ланцюжки, прямо не виміряно. & \cite{BiPoo1998} & виміряно \\
8 & Слід~\eqref{eq:threefactor}: дофамін, що надійшов через $1$--$2\,\text{с}$ після спільної активності, посилює зв'язок, пізніше --- ні (твердження~\ref{prop:threefactor}) & Зрізи стріатума, роздільна оптогенетична стимуляція глутаматних і дофамінових входів: шипик росте, лише якщо дофамін надходить через $0.3$--$2\,\text{с}$ після глутаматного входу; вікно задається швидким підйомом і розпадом цАМФ. У корі спарювання лишає окремі сліди для посилення і послаблення, які моноамінові рецептори перетворюють на довготривалі зміни у вікні порядку секунд. & \cite{Yagishita2014, He2015eligibility} & виміряно \\
9 & Вікна завдовжки в секунди бувають і без дофаміну: третій сигнал --- сильне збудження отримувача & Жива миша, поле CA1: одна плато-подія в отримувачі посилює входи, що надійшли за секунди до і після неї, і за один прохід створює поле місця. & \cite{Bittner2017} & виміряно \\
10 & Середній крок правила трьох множників іде за градієнтом винагороди~\eqref{eq:perturbation} (твердження~\ref{prop:gradient}) & Теорема про градієнт для навчання за випадковими відхиленнями; у розділі виведена для малого шуму. & \cite{Williams1992} & теорія \\
11 & Шум --- це пошук: мережа навчається на своїх випадкових відхиленнях & Молоді амадини: вимкнення ядра, що виходить із кола базальних гангліїв, різко і зворотно прибирає мінливість пісні. Дорослі амадини: завада подається на виконання складу з висотою по один бік від середньої, і птахи швидко зсувають висоту в інший бік --- навчаються за власним розкидом. & \cite{Olveczky2005, TumerBrainard2007} & виміряно \\
12 & Вибір з двох навчається при $\tau_e=1\,\text{с}$ і $\delta=R-\bar R$; не навчається при $\tau_e=50\ms$; без віднімання ваги ростуть без меж, а за великої швидкості навчання мережа може закріпити гірший вибір & \texttt{models/dofa\_learning.py}, $\eta=0.05$, $500$ спроб, $6$ зерен. $\tau_e=1\,\text{с}$: частка виборів $A$ $0.65$--$0.79\to0.96$--$1.00$, ваги $0.85$--$1.33$. $\tau_e=50\ms$: $0.38$--$0.55$, ваги не змінюються. $\delta=R$: вага переможця $860$--$1190$. $\delta=R$, $\eta=0.5$, $3$ зерна: одне закріпилося на $B$ ($0.04\to0$). Скрипт друкує всі чотири випадки; перерахунок збігся з розділом. & \texttt{models/\allowbreak dofa\_\allowbreak learning.py} & модель \\
13 & Час навчання за одним спільним сигналом росте пропорційно до кількості шумних нейронів (твердження~\ref{prop:broadcastcost}) & Криві навчання лінійної мережі: збурення нейронів повільніше за точний градієнт у стільки разів, скільки вихідних нейронів, збурення ваг --- ще в стільки разів, скільки входів. & \cite{Werfel2005} & теорія \\
14 & Виходи \nt{DOPA} ідуть насамперед до областей вибору дій & Повна реконструкція аксонів окремих нігростріатних \nt{DOPA}-нейронів щура: гілки однієї клітини покривають $0.45$--$5.7\,\%$ об'єму стріатума (у середньому $2.7\,\%$), поза стріатумом гілок майже немає. & \cite{Matsuda2009} & виміряно \\
15 & Кора навчається передбачати свій вхід, помилка рахується на місці & Огляд: у сенсорній корі є відповіді на розбіжність очікуваного і того входу, що надійшов. Що це загальне правило навчання кори, не доведено. & \cite{Keller2018} & гіпотеза \\
16 & Дофамін поводиться як помилка~\eqref{eq:ctd}: сплеск, перенесення на провісник, провал (рис.~\ref{fig:ctdcases}) & \nt{DOPA}-нейрони мавпи: сплеск на несподіваний сік; після навчання --- сплеск на провісник і немає відповіді на обіцяний сік; провал, коли сік не надійшов. Неперервна форма~\eqref{eq:ctd} --- теорія. & \cite{Schultz1997, Doya2000} & виміряно \\
17 & Схема, що обчислює $\dd V/\dd t$ і віднімає очікування & Миші, запис і оптогенетика у вентральній покришці: \nt{DOPA}-нейрони віднімають очікування від винагороди; сусідні \nt{GABA}-нейрони гальмують їх, коли винагорода очікується, і їхнє збудження чи гальмування змінює помилку. Як обчислюється похідна, не показано. & \cite{Eshel2015arithmetic} & гіпотеза \\
18 & \nt{DOPA}-нейрон --- ритмоводій; провали мілкіші за сплески & У зрізі \nt{DOPA}-нейрони розряджаються спонтанно і дуже регулярно, на повільному пейсмекерному потенціалі. У мавпи частота кількісно йде за додатною помилкою, а від'ємну передає лише слабко. & \cite{GraceOnn1989, Bayer2005} & виміряно \\
19 & Актор і критик (рис.~\ref{fig:actorcritic}) & Алгоритм --- з теорії навчання з підкріпленням. У людей (фМРТ) черевний стріатум поводиться як критик і під час вибору, і без нього, спинний --- лише коли треба вибирати дії. & \cite{SuttonBarto2018, ODoherty2004actor} & гіпотеза \\
20 & Знак $\delta$ у правилі перевернутий у нейронів із другою родиною рецепторів & Зрізи стріатума: у нейронів із D1-рецепторами спарювання дає посилення, якому потрібен D1; у нейронів із D2 --- послаблення, якому потрібен D2. & \cite{Shen2008} & виміряно \\
\bottomrule
\end{longtable}}

\section{Розділ~\ref{sec:dofaalt}. Інші теорії \nt{DOPA}}

Кожна з розширених картин спирається на свої прямі вимірювання дофаміну (розкид точок перевертання, відповіді на неприємне і нове, повільне зростання, відповіді на зміну смаку), але формули, які їх пояснюють, --- теорії, і щодо двох із них досліди поки суперечать одне одному.

{\footnotesize
\setlength{\tabcolsep}{3pt}\renewcommand{\arraystretch}{1.15}
\begin{longtable}{>{\raggedright}p{0.5cm}>{\raggedright}p{4.0cm}>{\raggedright}p{5.6cm}>{\raggedright}p{2.3cm}>{\raggedright\arraybackslash}p{1.7cm}}
\toprule
№ & Твердження & Дослід і що виміряно & Джерело & Статус \\
\midrule\endfirsthead
\toprule
№ & Твердження & Дослід і що виміряно & Джерело & Статус \\
\midrule\endhead
1 & Асиметричне правило~\eqref{eq:asymmetric} сходиться до точки~\eqref{eq:expectile}; при $\kappa=1/2$ це середнє, для краплі з імовірністю $1/2$ --- $V=\kappa$ (рис.~\ref{fig:expectiles}) & Точка~\eqref{eq:expectile} --- експектиль, розв'язок задачі найменших квадратів із різними вагами додатних і від'ємних відхилень; для прикладу розділу виведено в самому розділі. & \cite{NeweyPowell1987}; виведено в розділі & теорія \\
2 & У різних \nt{DOPA}-нейронів різні точки перевертання, і вони впорядковані за асиметрією (твердження~\ref{prop:expectile}) & Миші, оптогенетично мічені \nt{DOPA}-нейрони вентральної покришки, винагороди різного розміру: точка, де сплеск змінюється провалом, у різних нейронів різна; у нейронів із більшою асиметрією відповідей вона вища; за відповідями групи відновлюється форма розподілу винагород. Що це головна функція розкиду --- гіпотеза. & \cite{Dabney2020} & виміряно \\
3 & Частина \nt{DOPA}-нейронів відповідає на неприємне сплеском & Мавпи: одні \nt{DOPA}-нейрони збуджуються на провісник винагороди і гальмуються на провісник струменя повітря в око, інші збуджуються на обидва; групи розташовані в різних частинах середнього мозку. & \cite{Matsumoto2009, BrombergMartin2010} & виміряно \\
4 & Модуль помилки або новизна задає швидкість навчання & Прямого досліду, де сигнал важливості в \nt{DOPA}-нейронів змінює швидкість навчання, у цитованих роботах немає; огляд пропонує роль сигналу «увага, важливо». & \cite{BrombergMartin2010} & гіпотеза \\
5 & Окремий провід для осі небезпеки & Миші: \nt{DOPA}-нейрони, що йдуть до заднього кінця стріатума, відповідають на нові й сильні подразники, а не на винагороду; їхнє руйнування послаблює уникання нового предмета. & \cite{Menegas2018} & виміряно \\
6 & Оптимальний час дії $\tau_a^*=\sqrt{C/\bar R}$~\eqref{eq:vigor} & Мінімум суми $C/\tau_a+\bar R\tau_a$; виведено в розділі. У вихідній моделі швидкість дій задає ціна часу, що дорівнює середній винагороді. & \cite{Niv2007}; виведено в розділі & теорія \\
7 & Повільний рівень дофаміну несе $\bar R$ і задає швидкість дій (твердження~\ref{prop:multiplex}) & Щури, мікродіаліз і вольтамперометрія в прилеглому ядрі: рівень дофаміну від хвилини до хвилини йде за частотою винагород і за швидкістю дій. У людей L-DOPA посилює залежність часу реакції від середньої частоти винагород. Що повільний рівень дорівнює саме $\bar R$, не доведено. & \cite{Hamid2016, Beierholm2013vigor} & гіпотеза \\
8 & Повільний рівень складається з двох джерел: викид змінюється біля виходів без зміни частоти імпульсів, але повільне зростання буває і в самій частоті & Щури, одна задача: викид у прилеглому ядрі йде за змінним очікуванням винагороди, а частота імпульсів упізнаних \nt{DOPA}-нейронів вентральної покришки --- ні. Миші у віртуальному коридорі (рядок~10): плавне зростання під час наближення до винагороди є і в імпульсах упізнаних \nt{DOPA}-нейронів. & \cite{Mohebi2019, Kim2020} & виміряно \\
9 & Дофамін плавно росте, поки тварина наближається до далекої винагороди & Щури в лабіринті, вольтамперометрія в стріатумі: концентрація дофаміну наростає за секунди в міру наближення до мети, крутизна залежить від відстані і розміру винагороди. & \cite{Howe2013ramp, Hamid2016} & виміряно \\
10 & Зростання --- помилка $\delta$ при уточненні оцінки~\eqref{eq:ramp}; стрибок у коридорі дає сплеск (рис.~\ref{fig:ramp}) & Формула і число $\delta=0.75\,V$ на секунду --- виведено в розділі. Дослід: миші у віртуальному коридорі, миттєве перенесення ближче до мети; імпульси тіл, кальцій тіл і виходів та концентрація в стріатумі відповідають як помилка, а не як очікування $V$. Яке з двох пояснень правильне на всіх виходах, обговорюється. & \cite{Kim2020} & виміряно \\
11 & Погляд назад: дофамін повідомляє міру причинного зв'язку~\eqref{eq:retro} & Миші, викид дофаміну в прилеглому ядрі в дослідах, де ця теорія і помилка~\eqref{eq:ctd} передбачають різне: у низці дослідів викид ішов за першою. & \cite{Jeong2022} & гіпотеза \\
12 & Відповідь дофаміну під час навчання поступово зсувається від винагороди до провісника & Миші, сигнал виходів \nt{DOPA}-нейронів у черевному стріатумі в ході навчання: сплеск поступово переходить з моменту винагороди на момент провісника, як вимагає навчання за~\eqref{eq:ctd}. & \cite{Amo2022} & виміряно \\
13 & \nt{DOPA}-нейрони відповідають на зміну смаку винагороди за тієї самої цінності & Щури, запис нейронів вентральної покришки: заміна соку іншим, рівним за цінністю, спричиняє відповідь, хоча помилка~\eqref{eq:ctd} дорівнює нулю. & \cite{Takahashi2017} & виміряно \\
14 & Штучні сплески вчать зв'язку двох нейтральних подразників & Щури, дослід із попереднім поєднанням двох подразників без винагороди: оптогенетичне ввімкнення \nt{DOPA}-нейронів на переході від першого до другого створює зв'язок, вимкнення заважає його вивчити. & \cite{Sharpe2017} & виміряно \\
15 & Вектор помилок за ознаками~\eqref{eq:vectortd} & Теорія помилки передбачення ознак; прямої перевірки векторної форми немає. & \cite{Gardner2018} & гіпотеза \\
16 & Різні \nt{DOPA}-нейрони в одній задачі несуть різні величини & Миші у віртуальному коридорі, двофотонна зйомка \nt{DOPA}-нейронів: одні пов'язані з винагородою, інші зі швидкістю і поворотом, треті з провісниками і точністю вибору. & \cite{Engelhard2019} & виміряно \\
17 & Джгут замість проводу (твердження~\ref{prop:bundle}) & Зведення рядків~2--16: кожне розкладання виміряно окремо; єдиної картини, що всі вони --- частини одного сигналу, дослідом не встановлено. & --- & гіпотеза \\
\bottomrule
\end{longtable}}

\section{Розділ~\ref{sec:nora}. \nt{NORA}}

Будову \nt{NORA}-системи, її два режими, зв'язок із зіницею (не лише через \nt{NORA}), зростання відношення сигналу до фону і перевернуте U в поведінці виміряно прямо; мультиплікативне підсилення $g$ --- модель; те, що підсилення перемикає режим вибору і працює як обернена температура, --- висновки розділу; виграш від підлаштування підсилення --- розрахунок за моделлю книги; «\nt{NORA} --- ручка пошуку» і «\nt{NORA} --- переривання» --- гіпотези.

{\footnotesize
\setlength{\tabcolsep}{3pt}\renewcommand{\arraystretch}{1.15}
\begin{longtable}{>{\raggedright}p{0.5cm}>{\raggedright}p{4.0cm}>{\raggedright}p{5.6cm}>{\raggedright}p{2.3cm}>{\raggedright\arraybackslash}p{1.7cm}}
\toprule
№ & Твердження & Дослід і що виміряно & Джерело & Статус \\
\midrule\endfirsthead
\toprule
№ & Твердження & Дослід і що виміряно & Джерело & Статус \\
\midrule\endhead
1 & \nt{NORA}-нейронів у людини кілька десятків тисяч, майже всі в одній групі & Серійні зрізи стовбура мозку людини, забарвлення на тирозингідроксилазу: $53\,900$ клітин у блакитній плямі і ще $6\,260$ у сусідньому під'ядрі. & \cite{Baker1989LC} & виміряно \\
2 & Виходи розходяться майже по всьому мозку, але частини групи частково обслуговують різні області & Вірусно-генетичне мічення в мишей: \nt{NORA}-нейрони блакитної плями отримують входи з широких областей мозку і надсилають виходи в широкі області головного і спинного мозку. У щурів: група, що йде до мигдалеподібного тіла, посилює вивчення страху, а окрема група, що йде до префронтальної кори, --- його згасання. & \cite{Schwarz2015LC, Uematsu2017LC, Poe2020} & виміряно \\
3 & Фоновий режим: рівні імпульси з частотою від часток до кількох на секунду, що змінюється з рівнем неспання & Щури, запис нейронів блакитної плями в циклі сну і неспання: частота найбільша в неспанні, нижча в повільному сні, майже нуль у швидкому; зміни частоти випереджають зміну стану. & \cite{AstonJonesBloom1981} & виміряно \\
4 & Уривчастий режим: коротка пачка на важливу для задачі подію, приблизно в момент рішення & Мавпи, задача знайти рідкісну ціль: усі нейрони блакитної плями відповідають пачкою лише на ціль, через $\sim90\ms$ після неї і за $\sim200\ms$ до відповіді рукою; за поганої роботи пачки слабші. У задачі з вибором із двох пачка точніше прив'язана до відповіді, ніж до стимулу, і є як за правильних, так і за помилкових відповідей. & \cite{AstonJones1994vigilance, Clayton2004LC} & виміряно \\
5 & Діаметр зіниці --- неточна контрольна точка \nt{NORA}: він іде і за \nt{NORA}, і за \nt{ACET}, і за іншими областями & Мавпи: імпульси блакитної плями, аж до окремих імпульсів однієї клітини, передбачають розширення зіниці; схожий, але пізніший зв'язок є в горбках і поясній корі. Миші: коливання зіниці йдуть за активністю і норадреналінових, і холінергічних виходів у корі. & \cite{Joshi2016, Reimer2016} & виміряно \\
6 & Норадреналін пригнічує фонову активність сильніше за викликану; мультиплікативне підсилення~\eqref{eq:gain} --- модель, що передає цей ефект & Мавпи в стані неспання, слухова кора, іонофорез норадреналіну: фонова активність нейронів пригнічується сильніше, ніж відповідь на голоси родичів, --- відношення сигналу до шуму зростає. Мультиплікативну сигмоїду~\eqref{eq:gain} взято з мережевої моделі; буквального множення крутизни не виміряно. & \cite{Foote1975NE, ServanSchreiber1990} & виміряно; $g$ --- модель \\
7 & Підсилення підвищує $d'$ лише проти шуму після підсилювача~\eqref{eq:gainsnr} (твердження~\ref{prop:gainnoise}) & Виведено в розділі; у мережевій моделі підсилення покращує виявлення сигналу. Досліду, який відокремлював би шум до і після підсилювача і перевіряв би цю різницю, немає. & виведено в розділі; \cite{ServanSchreiber1990} & теорія \\
8 & Межа $g\beta=1$ перемикає мережу вибору з «розмірковує» на «вирішила»~\eqref{eq:wtagain} (твердження~\ref{prop:gainwta}, рис.~\ref{fig:pitchfork}) & Лінійний аналіз двох груп, що гальмують одна одну; виведено в розділі. Прямих вимірювань цього порога в мозку немає. & виведено в розділі & теорія \\
9 & Коефіцієнт підсилення --- обернена температура вибору~\eqref{eq:softmax} & За логістичного шуму після підсилювача ймовірність вибору --- softmax з $T=\sigma/g$; виведено в розділі. & виведено в розділі & теорія \\
10 & \nt{NORA} задає, скільки шукати: висока фонова активність --- мале діюче $g$ (пошук), уривчаста пачка в момент рішення --- велике (рішення) & Люди: перед вибором навмання базова зіниця ширша, ніж перед вибором найкращого відомого; у людей із ширшою зіницею схильність шукати вища. Щури: перехід у режим випадкового вибору керується входом \nt{NORA} у передню поясну кору. Що пачка підвищує діюче $g$, прямо не виміряно. & \cite{JepmaNieuwenhuis2011, Tervo2014stochastic, AstonJones2005} & гіпотеза \\
11 & Винагорода залежить від $g$ як перевернуте U; найкраще $g$ менше у світі, що швидко змінюється (твердження~\ref{prop:explore}, рис.~\ref{fig:invertedU}) & \texttt{models/nora\_explore.py}, $60$ прогонів по $4000$ спроб. Зміна раз на $1000$ спроб: найкраще $g=16$, частка $0.976$; раз на $20$: $g=8$, частка $0.886$; при $g=0.5$ --- $0.77$. Перерахунок збігся з розділом. & \texttt{models/\allowbreak nora\_\allowbreak explore.py} & модель \\
12 & Перевернуте U в поведінці: найкраще за помірної фонової активності з чіткими пачками & Мавпи, та сама задача, що в рядку~4: у періоди доброї роботи фонова частота помірна, а пачки на ціль чіткі; за високої фонової частоти пачок немає і хибних відповідей більше. Зміни фонової і викликаної активності йдуть за коливаннями якості роботи. & \cite{Usher1999LC, AstonJones2005} & виміряно \\
13 & \nt{NORA} повідомляє неочікувану невизначеність, \nt{ACET} --- очікувану & Теорія байєсівського висновку з двома видами невизначеності; прямого досліду, що розділяє ці сигнали за медіаторами, у цитованих роботах немає. & \cite{YuDayan2005} & гіпотеза \\
14 & Після різкої зміни люди навчаються швидше, і зіниця розширюється & Люди, задача передбачення з раптовими змінами: короткі зміни зіниці йдуть за оцінкою ймовірності зміни і разом із базовою зіницею передбачають, наскільки сильно нові дані змінюють оцінку (швидкість навчання); зміна зіниці, не пов'язана зі світлом і задачею, змінює цю швидкість. Що зіниця тут показує саме \nt{NORA}, --- висновок за непрямими даними рядка~5. & \cite{Nassar2012} & виміряно \\
15 & Уривчаста пачка працює як переривання: мережа скидає стан & Прямого досліду, де пачка \nt{NORA} скидає стан мережі, немає; виміряно лише пачки на важливі події (рядок~4). & \cite{BouretSara2005} & гіпотеза \\
16 & Підлаштування $g$ або швидкості навчання за несподіваністю майже не краще за найкращі сталі налаштування & \texttt{models/nora\_explore.py}, світ з епохами по $1000$ спроб (зміна раз на $1000$ і раз на $20$). Найкраще стале $g=8$: $0.925$; підлаштування $g$: до $0.927$; підлаштування швидкості навчання: до $0.926$; стала швидкість $0.4$: $0.928$. Перерахунок збігся з розділом. & \texttt{models/\allowbreak nora\_\allowbreak explore.py} & модель \\
\bottomrule
\end{longtable}}

\section{Розділ~\ref{sec:acet}. Додаємо \nt{ACET}}

Три дії ацетилхоліну виміряно на зрізах і в живому мозку, конфлікт запису і читання показано на моделі книги, а те, що \nt{ACET} слугує лінією дозволу запису і повідомляє очікувану невизначеність, --- гіпотеза з частковою опорою на дослід.

{\footnotesize
\setlength{\tabcolsep}{3pt}\renewcommand{\arraystretch}{1.15}
\begin{longtable}{>{\raggedright}p{0.5cm}>{\raggedright}p{4.0cm}>{\raggedright}p{5.6cm}>{\raggedright}p{2.3cm}>{\raggedright\arraybackslash}p{1.7cm}}
\toprule
№ & Твердження & Дослід і що виміряно & Джерело & Статус \\
\midrule\endfirsthead
\toprule
№ & Твердження & Дослід і що виміряно & Джерело & Статус \\
\midrule\endhead
1 & \nt{ACET}-нейронів мозку небагато, вони зібрані в кількох групах, а виходи розходяться по всій корі: широкомовний канал & Забарвлення антитілами до ферменту синтезу ацетилхоліну і ретроградна мітка в щура: кора, гіпокамп, мигдалеподібне тіло, нюхова цибулина і таламус отримують ацетилхолін переважно від шести компактних груп клітин базального переднього мозку і стовбура & \cite{Mesulam1983} & виміряно \\
2 & Рівень ацетилхоліну в корі високий у неспанні і низький у повільному сні & Мікродіаліз у корі і гіпокампі вільно рухливих котів із записом ЕЕГ: викид ацетилхоліну в спокійному неспанні на $\sim100\%$, в активному --- на $\sim175\%$ вищий, ніж у повільному сні; у швидкому сні в корі --- як у неспанні & \cite{Marrosu1995,Hasselmo1999} & виміряно \\
3 & Зміст сигналу задає рецептор: ацетилхолін має швидкі рецептори-канали і повільні, що запускають реакції в клітині (твердження~\ref{prop:receptor}) & Огляд вимірювань: дія ацетилхоліну на збудливість, передачу і пластичність залежить від місця викиду, підтипу рецептора (нікотинові --- іонні канали, мускаринові --- через G-білки) і клітини-отримувача & \cite{Picciotto2012} & виміряно \\
4 & Перша дія: послабити внутрішні зв'язки, майже не зачіпаючи входів ззовні (множник $1-a$ у~\eqref{eq:acetnet}) & Зрізи нюхової кори щура: за $100$\,мкМ агоністів ацетилхоліну потенціали внутрішніх зв'язків (шар Ib) падають більш ніж на $60\%$, зовнішнього входу (шар Ia) --- менш ніж на $15\%$, у тій самій клітині; пригнічення пресинаптичне. Зрізи гіпокампа (CA1): $89\%$ проти $40\%$ для внутрішнього і вхідного шляхів & \cite{HasselmoBower1992,HasselmoSchnell1994} & виміряно \\
5 & Друга дія: посилити вхід (множник $\frac{1+a}{2}$) & Зрізи неокортексу: нікотинові рецептори посилюють синапси входу з таламуса і не діють на внутрішньокіркові; мускаринові послаблюють обидва типи. Збудливість нейронів ацетилхолін підвищує (огляд) & \cite{Gil1997,Hasselmo2006} & виміряно \\
6 & Третя дія: полегшити зміну ваг (швидкість $\eta a$) & Щур: електрична стимуляція базального ядра (джерела ацетилхоліну для кори) в парі з тоном спричиняє в дорослої тварини сильну перебудову карти слухової кори під звук, що пред'являється & \cite{KilgardMerzenich1998,Hasselmo2006} & виміряно \\
7 & Правило Гебба для рідкісних образів у другому рівнянні~\eqref{eq:acetnet} дає асоціативну пам'ять, що дописує образ за частиною & Розрахунок ємності мережі з розрідженими образами і коваріаційним правилом: за малої частки активних нейронів $p$ мережа зберігає помітно більше образів, ніж при $p=1/2$ & \cite{Tsodyks1988} & теорія \\
8 & Запис за низького $a$ псує пам'ять: мережа записує згадане, а не побачене; при $a=0.1$ псування не слабше, ніж при $a=0$ & \texttt{models/\allowbreak acet\_memory.py}: $200$ нейронів, п'ять образів по $20$, новий образ ділить з одним із них $10$ нейронів, $10$ прогонів. При $a=0$ новий образ не запам'ятано, старий тягне $5.9$ чужого нейрона з $10$; при $a=0.1$ новий образ згадується ($0.97$), але обидва зливаються: новий тягне $7.3$ чужого нейрона, старий --- $7.9$; з $a=0.2$ --- менше за один & виведено в розділі & модель \\
9 & Твердження~\ref{prop:writeread}: немає рівня $a$, за якого однаково добрі запис і читання (рис.~\ref{fig:acetsim}) & Той самий скрипт, швидкість навчання $\eta a$: новий образ запам'ятовується за одне пред'явлення цілком при $a\ge0.9$, на $0.8$ при $a=0.75$, не запам'ятовується при $a\le0.5$. Читання за половиною образу: $1.0$ при $a\le0.2$, $0.96$ при $a=0.3$, $0.04$ при $a=0.4$ & виведено в розділі & модель \\
10 & У мозку один широкомовний провід --- \nt{ACET} --- перемикає запис і читання & Ідею взято з моделей, побудованих за вимірюваннями на зрізах. Найпряміший дослід --- у людини: блокатор мускаринових рецепторів скополамін погіршує заучування нових пар слів, найсильніше тих, що перетинаються зі старими, але не згадування вже вивчених пар. Прямого вимірювання двох режимів мережі немає & \cite{HasselmoAndersonBower1992,Atri2004} & гіпотеза \\
11 & Фільтр~\eqref{eq:kalman} оптимальний; вага входу і швидкість навчання --- одна величина $P/R$, в усталеному режимі $P=\sqrt{qR}$ і пам'ять $\sqrt{R/q}$ (твердження~\ref{prop:kalman}) & Досліду не потребує: оптимальність фільтра Калмана--Б'юсі --- класичний результат теорії оцінювання; усталений режим виведено в розділі & \cite{KalmanBucy1961} & теорія \\
12 & \nt{ACET} повідомляє мережі величину, подібну до $P$, --- очікувану невизначеність & Запропоновано в імовірнісній моделі уваги. Прямого вимірювання того, що рівень ацетилхоліну йде за невизначеністю оцінки, немає & \cite{YuDayan2005} & гіпотеза \\
13 & Частина \nt{ACET}-нейронів відповідає на винагороду і покарання за два десятки мілісекунд, тим сильніше, чим несподіваніше підкріплення & Миша, оптогенетично впізнані холінергічні нейрони базального переднього мозку: відповідь на винагороду і покарання через $18\pm3$\,мс, величина зростає з несподіваністю підкріплення, відповіді схожі в нейронів двох ядер & \cite{Hangya2015} & виміряно \\
14 & \nt{NORA} помічає неочікувану зміну світу, \nt{ACET} тримає мережу в режимі запису & Розподіл праці запропоновано в моделі. Непрямо: у людини діаметр зіниці, пов'язаний із \nt{NORA}, відстежує зміни і швидкість навчання. Прямої перевірки пари \nt{NORA}--\nt{ACET} немає & \cite{YuDayan2005,Nassar2012} & гіпотеза \\
15 & У повільному сні мережа в режимі читання програє записане вдень, і ці повтори переписують його в довготривалу пам'ять & Записи ансамблів нейронів гіпокампа щура: клітини, що працювали разом удень, частіше спрацьовують разом у наступному повільному сні і в тому самому порядку --- виміряно. Для перезапису: пригнічення брижів гіпокампа після навчання погіршує просторову пам'ять, а подовження спонтанних брижів її покращує; що причина --- низький рівень \nt{ACET}, не доведено & \cite{WilsonMcNaughton1994,Skaggs1996,Girardeau2009,FernandezRuiz2019,Hasselmo1999} & гіпотеза \\
\bottomrule
\end{longtable}}

\section{Розділ~\ref{sec:synthesis}. Уся машина}

Розділ-зведення нових дослідів не наводить: кожен його рядок спирається на розділ, де твердження розібрано, і на ті самі джерела; твердо встановлено шину \nt{GLUT} і керування потоком через \nt{GABA}, а ролі широкомовних каналів та їхня спільна робота --- здебільшого гіпотеза.

{\footnotesize
\setlength{\tabcolsep}{3pt}\renewcommand{\arraystretch}{1.15}
\begin{longtable}{>{\raggedright}p{0.5cm}>{\raggedright}p{4.0cm}>{\raggedright}p{5.6cm}>{\raggedright}p{2.3cm}>{\raggedright\arraybackslash}p{1.7cm}}
\toprule
№ & Твердження & Дослід і що виміряно & Джерело & Статус \\
\midrule\endfirsthead
\toprule
№ & Твердження & Дослід і що виміряно & Джерело & Статус \\
\midrule\endhead
1 & На $8.6\cdot10^{10}$ нейронів припадає лише чотири види керувальних сигналів~\eqref{eq:machine} & Підрахунок ядер клітин у гомогенаті мозку (ізотропний фракціонатор): у дорослого чоловіка $86.1\pm8.1$ млрд нейронів & \cite{Azevedo2009} & виміряно \\
2 & Твердження~\ref{prop:architecture}, перші два шари: дані йдуть адресною шиною \nt{GLUT}, знак зв'язку задає відправник, потік спрямовують і нормують \nt{GABA}-нейрони (розділи~\ref{sec:neurontypes}, \ref{sec:gaba}) & Один головний медіатор на нейрон (огляд вимірювань); пряме гальмування звужує вікно, в якому пірамідний нейрон додає входи; ділення на загальну активність виміряно в багатьох ділянках & \cite{Strata1999,Pouille2001,Carandini2012} & виміряно \\
3 & Елементи ненадійні: синапс губить більшість імпульсів (табл.~\ref{tab:computer}, розділ~\ref{sec:code}) & Зрізи гіпокампа, мінімальна стимуляція: понад половина імпульсів не викликає відповіді в CA1; збої порога й проведення аксоном виключено, причина --- імовірнісне вивільнення медіатора & \cite{AllenStevens1994} & виміряно \\
4 & Мозок працює на $\sim20$ ватах, у середньому близько імпульсу за секунду на нейрон (розділ~\ref{sec:code}); такої машини інженери поки не збудували & Оцінка~\eqref{eq:energy} з виміряних витрат: близько $10^9$ молекул АТФ на імпульс разом із роботою синапсів (кора гризунів); детальна оцінка для кори дає ще меншу частоту. Стан техніки --- за оглядом & \cite{Attwell2001,Lennie2003,MehonicKenyon2022} & теорія \\
5 & Навчання більшості синапсів --- добуток місцевого сліду й одного спільного числа $\delta$ (друге рівняння~\eqref{eq:machine}, розділ~\ref{sec:dofa}) & \nt{DOPA}-нейрони мавпи відповідають на помилку передбачення винагороди; у зрізах смугастого тіла дофамін збільшує шипик, лише якщо надходить через $0.3$--$2$\,с після глутаматного входу --- слід виміряно & \cite{Schultz1997,Yagishita2014} & виміряно \\
6 & Окремий провід помилки до кожного виходу є лише в колі навчання рухів (розділ~\ref{sec:motorlearning}) & Мозочок: збіг сигналу ліаноподібного волокна з входом послаблює цей вхід; у мавпи складні імпульси клітин Пуркіньє несуть напрямок помилки руху й викликають поправку & \cite{Ito2001,Herzfeld2018} & виміряно \\
7 & Кожен широкомовний канал --- джгут із кількох ліній (твердження~\ref{prop:bundle}) & Для \nt{DOPA}: в одній задачі різні нейрони несуть винагороду, рух і ознаки стимулу; швидка помилка й повільний рівень дофаміну різні. Для \nt{SERO}, \nt{NORA}, \nt{ACET} такий розклад вивчено гірше & \cite{Engelhard2019,Hamid2016,Mohebi2019} & виміряно \\
8 & \nt{SERO}: регулятор рівня і, за гіпотезою, горизонт $\tau_\gamma$ (розділ~\ref{sec:serocomputer}) & Самогальмування: агоністи власних рецепторів 5-HT$_{1A}$ гальмують серотонінові нейрони шва --- виміряно. Горизонт запропоновано в теорії; найсильніший дослід --- оптогенетична активація цих нейронів у миші подовжує очікування відкладеної винагороди & \cite{Sprouse1987,Doya2002,Miyazaki2014} & гіпотеза \\
9 & \nt{NORA}: підсилення $g$ вибирає між пошуком і рішенням (розділ~\ref{sec:nora}) & Зростання відношення сигнал/шум: у мавпи в стані неспання норадреналін у слуховій корі пригнічує фонову активність сильніше, ніж відповідь на голоси родичів --- виміряно; мультиплікативне підсилення --- модель; роль у виборі між пошуком і рішенням --- теорія & \cite{Foote1975NE,ServanSchreiber1990,AstonJones2005} & гіпотеза \\
10 & \nt{ACET}: перемикає запис і читання (розділ~\ref{sec:acet}) & Три дії (послаблення внутрішніх зв'язків, підсилення входу, полегшення пластичності) виміряно; перемикання режимів опосередковано підтримане дослідом зі скополаміном у людини & \cite{HasselmoBower1992,Gil1997,Atri2004} & гіпотеза \\
11 & Формула~\eqref{eq:machine} описує машину цілком & Це зведення моделей розділів~\ref{sec:glutcomputer}--\ref{sec:acet}, кожна частина спирається на свій розділ; як ціле дослідом не перевірялася & виведення в розділі & гіпотеза \\
12 & Ручки працюють злагоджено без спільного керування: помилка, несподіванка, запис, повернення (рис.~\ref{fig:timeline}) & Досліду немає: схема якісна. Окремі ланки --- теорія розподілу праці \nt{NORA} і \nt{ACET} та зв'язок зіниці зі швидкістю навчання в людини & \cite{YuDayan2005,Nassar2012} & гіпотеза \\
13 & Навчання за одним спільним сигналом тим повільніше, що більше нейронів впливає на результат (твердження~\ref{prop:broadcastcost}) & Розрахунок кривих навчання лінійної мережі: навчання за збуренням виходів повільніше за градієнтний спуск у стільки разів, скільки виходів & \cite{Werfel2005} & теорія \\
14 & Як кора налаштовує багатошарові кола, невідомо; кандидати --- пачки імпульсів, обчислення в гілках нейрона, самонавчання передбаченням & Пачки як носій помилки --- модель; повторити відповідь детальної моделі нейрона кори змогла лише мережа з $5$--$8$ шарів --- модель; кальцієві імпульси в гілках нейронів кори людини, що дають виключне АБО, --- виміряно в зрізах; самонавчання --- огляд свідчень & \cite{Payeur2021,Beniaguev2021,Gidon2020,Keller2018} & гіпотеза \\
\bottomrule
\end{longtable}}

\section{Розділ~\ref{sec:sensors}. Входи машини}

Будову датчиків виміряно прямо --- записами струму поодиноких клітин, коливань завитки й підрахунком клітин сітківки; межа чутливості й формула дробового шуму випливають із фізики, а те, що коди датчиків підлаштовано під найбільшу інформацію, --- гіпотеза з сильною опорою.

{\footnotesize
\setlength{\tabcolsep}{3pt}\renewcommand{\arraystretch}{1.15}
\begin{longtable}{>{\raggedright}p{0.5cm}>{\raggedright}p{4.0cm}>{\raggedright}p{5.6cm}>{\raggedright}p{2.3cm}>{\raggedright\arraybackslash}p{1.7cm}}
\toprule
№ & Твердження & Дослід і що виміряно & Джерело & Статус \\
\midrule\endfirsthead
\toprule
№ & Твердження & Дослід і що виміряно & Джерело & Статус \\
\midrule\endhead
1 & Датчик --- перетворювач: рецептор дає струм, а вихід чутливих клітин ока й вуха --- глутамат на шину \nt{GLUT}; перші клітини імпульсів не видають (рис.~\ref{fig:transducer}) & Палички й колбочки черепахи виділяють збуджувальну амінокислоту: її вловлюють глутаматні канали на відокремленому клаптику мембрани. Внутрішні волоскові клітини завитки щура: струми в закінченнях нервових волокон ідуть через глутаматні AMPA-рецептори, частота вивільнень зростає з плавною деполяризацією клітини до $150$\,Гц & \cite{CopenhagenJahr1989,GlowatzkiFuchs2002} & виміряно \\
2 & Датчик світла в темряві відповідає на один фотон струмом близько пікоампера & Всмоктувальний електрод на зовнішньому сегменті палички ропухи: відповіді на тьмяні спалахи квантовані, подія $\approx1$\,пА ($5\%$ темнового струму) з розкидом $0.2$\,пА, ефективність $0.5$. Людина повідомляє про влучання одного фотона частіше, ніж випадково & \cite{Baylor1979,Tinsley2016} & виміряно \\
3 & Датчик звуку помічає зміщення менше за нанометр & Метод Мессбауера, основна мембрана завитки морської свинки в місці $16$--$19$\,кГц: на порозі відповіді слухового нерва швидкість мембрани близько $0.04$\,мм/с, тобто зміщення $\approx0.35$\,нм (наш перерахунок). Вимірювали мембрану, а не пучок датчика & \cite{Sellick1982,Hudspeth1989} & виміряно \\
4 & Точність у темряві обмежена дробовим шумом фотонів~\eqref{eq:photonnoise}; за слабкого світла накопичення довше & Формула випливає з пуассонівської статистики. У паличці шум за тьмяного світла збігається із сумою незалежних квантових подій; на яскравому тлі відповідь на спалах менша й коротша. Число «десяті частки секунди» тут окремо не перевірено & \cite{Baylor1979} & теорія \\
5 & Діапазон входу --- десять порядків для світла й дванадцять для звуку, а частота імпульсів --- менше за три порядки & Для звуку: відповідь основної мембрани на характеристичній частоті зростає з крутістю до $0.2$\,дБ/дБ у смузі $40$--$80$\,дБ, стиснення видно й вище $100$\,дБ. Самі числа порядків --- стандартні оцінки, у розділі без джерела & \cite{Ruggero1997} & виміряно \\
6 & Відповідь датчика~\eqref{eq:nakarushton}, а $\sigma$ іде за середньою яскравістю, зсуваючи криву логарифмічною віссю (рис.~\ref{fig:adaptation}) & Формулу вперше підібрано до S-потенціалів сітківки риби. Колбочки черепахи: відповіді підкоряються $V=I V_m/(I+\sigma)$ і охоплюють $\sim3.5$ порядку яскравості; після $2$--$3$ хвилин тла крива зсувається по горизонталі, зберігаючи ширину. Зв'язок із діленням на загальну активність --- огляд & \cite{NakaRushton1966,NormannPerlman1979,Carandini2012} & виміряно \\
7 & Датчики передають зміни, а їхні коди підлаштовано під найбільшу інформацію на імпульс (твердження~\ref{prop:sensors}) & Принцип запропоновано теоретично. Найсильніше свідчення: у мухи характеристика «контраст $\to$ відповідь» нейронів першого шару збігається з кумулятивним розподілом контрастів природних сцен --- так досягається найбільша інформація & \cite{Barlow1961,Laughlin1981} & гіпотеза \\
8 & Датчики ввімкнення й вимкнення --- двопровідний код; подієва камера повторює його в кремнії & Огляд дослідів: канали ввімкнення й вимкнення сітківки розділяються вже на першому синапсі й вимикаються окремо. Камера: $128\times128$ пікселів без кадрів, діапазон $120$\,дБ, затримка $15$\,мкс & \cite{Schiller1992,Lichtsteiner2008,Gallego2022} & виміряно \\
9 & В оці $\sim10^8$ датчиків світла й $\sim10^6$ вихідних нейронів: стиснення в $\sim100$ разів & Підрахунок на цілих препаратах сітківки людини: у середньому $92$ млн паличок і $4.6$ млн колбочок; вихідних (гангліозних) клітин від $0.7$ до $1.5$ млн у різних очах & \cite{Curcio1990,CurcioAllen1990} & виміряно \\
10 & У миші вихідних каналів ока понад тридцять & Миша: двофотонна кальцієва зйомка понад $11\,000$ вихідних клітин і кластеризація відповідей --- помітно більше за $30$ функціональних типів. Для людини число не виміряно & \cite{Baden2016} & виміряно \\
11 & Око морської свинки передає $\sim875\,000$ біт/с; перерахунок на око людини дає $\sim10^7$ біт/с & Морська свинка, багатоелектродна матриця, рух у природних сценах: $6$--$13$ біт/с на клітину, $\sim875\,000$ біт/с на $\sim10^5$ вихідних клітин. Для людини ($\sim10^6$ клітин) $\sim10^7$ біт/с --- перерахунок авторів & \cite{Koch2006} & виміряно \\
12 & Вухо --- банк смугових фільтрів із майже логарифмічною картою частот (рис.~\ref{fig:filterbank}) & Місце найбільшого коливання основної мембрани залежить від частоти; функція «частота--місце» майже експоненційна й однією формою описує завитки людини та живих тварин & \cite{Greenwood1990,Robles2001} & виміряно \\
13 & Резонатори активні: частина датчиків підсилює слабкі коливання в тисячі разів за амплітудою ($66$--$76$\,дБ), зі зростанням гучності підсилення падає & Лазерна віброметрія основи завитки шиншили: за слабких тонів на характеристичній частоті підсилення $66$--$76$\,дБ відносно стремінця; смерть знижує чутливість на $60$--$81$\,дБ (у $10^3$--$10^4$ разів за амплітудою) і прибирає стиснення. Джерело сили --- зовнішні волоскові клітини & \cite{Ruggero1997,Robles2001} & виміряно \\
14 & Підсилювач вуха стоїть біля межі самозбудження & Розрахунок: біля точки біфуркації Хопфа неминучі стиснення діапазону, гостре налаштування й комбінаційні тони --- усі відомі нелінійності слуху. Що завитка налаштована саме так, прямо не доведено & \cite{Eguiluz2000} & гіпотеза \\
15 & На низьких частотах імпульси йдуть у фазі зі звуком (у морської свинки прив'язка слабшає від $\sim600$\,Гц і помітна до $\sim3.5$\,кГц), і за ними знаходять напрямок; на високих лишається код місцем & Слуховий нерв морської свинки: прив'язка до фази починає падати близько $600$\,Гц і зникає вище $3.5$\,кГц; у кішки й мавпи межа приблизно на октаву вища. Різниця часу приходу у два вуха --- огляд & \cite{PalmerRussell1986,Grothe2010} & виміряно \\
16 & Інші датчики (табл.~\ref{tab:sensors}): у нюху $\sim400$ типів рецепторів, по одному на нейрон, і комбінаторний код; смак частково через АТФ і серотонін; дотик --- датчики, що звикають швидко й повільно; тепло й холод --- окремі & Миша: один рецептор упізнає багато запахів, один запах --- багато рецепторів, різні запахи --- різні набори. Без рецепторів АТФ нерв смаку не відповідає; смакові бруньки виділяють серотонін. Запис поодиноких волокон шкіри руки людини: чотири типи за швидкістю звикання. Холодовий канал відмінний від теплових & \cite{Mombaerts2004,Malnic1999,Finger2005,Huang2005,JohanssonVallbo1979,McKemy2002} & виміряно \\
\bottomrule
\end{longtable}}

\section{Розділ~\ref{sec:recognition}. Розпізнавання}

Швидкість розпізнавання, будову перших ступенів, лінійне читання відповіді й навчання на неперервності часу виміряно на людині й мавпі; зведення всього ланцюжка до двох операцій і навчання на відео перевірено на моделях книги, а пояснення ілюзій --- від добре обґрунтованих до спірних.

{\footnotesize
\setlength{\tabcolsep}{3pt}\renewcommand{\arraystretch}{1.15}
\begin{longtable}{>{\raggedright}p{0.5cm}>{\raggedright}p{4.0cm}>{\raggedright}p{5.6cm}>{\raggedright}p{2.3cm}>{\raggedright\arraybackslash}p{1.7cm}}
\toprule
№ & Твердження & Дослід і що виміряно & Джерело & Статус \\
\midrule\endfirsthead
\toprule
№ & Твердження & Дослід і що виміряно & Джерело & Статус \\
\midrule\endhead
1 & Порівняння зі зразком за пікселями при зсуві вгадує не краще за монетку; пара ступенів~\eqref{eq:scstage} упізнає фігуру в будь-якому місці & \texttt{models/\allowbreak recognition\_models.py}, «сітківка» з $24$ точок, $4000$ проб: зразок --- $0.49$, $0.51$, $0.52$ за частки перевернутих точок $0$, $0.1$, $0.2$; пара $S$ і $C$ --- $1.00$, $0.86$, $0.70$ & виведення в розділі & модель \\
2 & Тварину на картинці впізнають за $\sim150$\,мс, один прохід десятком ступенів по $10$--$20$\,мс & Людина, задача «є тварина чи ні» на фотографіях, показаних на $20$\,мс: викликані потенціали для двох відповідей розходяться через $\sim150$\,мс. Мавпа: час першої відповіді зростає від ступеня до ступеня (V1 раніше за всіх, далі V2 і V4). Кількість ступенів і «нуль або один імпульс на ступінь» --- висновок із цих чисел & \cite{Thorpe1996,Schmolesky1998} & виміряно \\
3 & Ступінь $S$ --- І за частинами, ступінь $C$ --- найбільше за положеннями (твердження~\ref{prop:conveyor}) & Кішка, первинна зорова кора: прості клітини відповідають на край певної орієнтації в певному місці, складні --- на ту саму орієнтацію в більшій ділянці. Внутрішньоклітинний запис складних клітин: відповідь на дві смужки в багатьох клітин близька до більшої з двох відповідей, у середньому найближча до операції max. Найбільше через взаємне гальмування --- твердження~\ref{prop:wta}, ділення на загальну активність --- огляд & \cite{HubelWiesel1962,Lampl2004,Carandini2012} & виміряно \\
4 & Вище ланцюжком поєднання більші й складніші, відповідь дедалі менше залежить від положення & Мавпа, спрощення стимулів до «критичної ознаки»: клітини V4 і задньої нижньоскроневої кори відповідають на складні поєднання форми, кольору й текстури за малого поля; у передній нижньоскроневій корі поле більше. Чергування $S$ і $C$ у моделі відтворює цю картину & \cite{KobatakeTanaka1994,Riesenhuber1999} & виміряно \\
5 & Згорткові мережі повторюють це чергування й найкраще передбачають відповіді верхніх ступенів; предмет читається лінійно за частотами групи нейронів & Мережа, навчена розпізнавати, без підгонки до мозку: верхній шар передбачає відповіді нейронів нижньоскроневої кори мавпи, середні --- відповіді V4. Лінійний класифікатор за активністю $\sim100$ випадково вибраних клітин у вікнах від $12.5$\,мс упізнає предмет і категорію за різних положення й розміру & \cite{Fukushima1980,Yamins2014,Hung2005} & виміряно \\
6 & Правило Гебба зі слідом~\eqref{eq:traceinvariance} навчає інваріантності на відео без підписів, якщо слід не коротший за $\sim20$\,мс & Ідея повільних ознак і правила зі слідом --- теорія. \texttt{models/\allowbreak recognition\_models.py}, два $C$-нейрони, $16$ входів, $10$ прогонів: пауза між предметами $0.3$\,с. При $\tau_{\text{сл}}=5$\,мс збіг із фігурою $0.52$ у середньому, при $10$\,мс --- $0.56$; при $20$, $50$ і $200$\,мс --- $1.00$ в усіх прогонах (при $30$\,мс один прогін із десяти помиляється) & \cite{Foldiak1991,WiskottSejnowski2002} & модель \\
7 & У мозку інваріантність навчається на неперервності часу & Мавпа: предмет непомітно підміняли під час стрибка ока в певне місце поля; позиційна інваріантність нейронів нижньоскроневої кори змінювалася відповідно до підміни, значуще вже за годину. Що це головне правило навчання зору --- гіпотеза & \cite{LiDiCarlo2008} & виміряно \\
8 & Рух: детектори напрямку, далі та сама пара І/АБО за великими ділянками & Детектори напрямку в сітківці кроля; в ділянці MT мавпи всі $110$ записаних нейронів вибіркові до напрямку, відповідь на рух сильніша, ніж у V1 & \cite{Barlow1965,Albright1984} & виміряно \\
9 & Верхні ступені надсилають униз передбачення, угору йде помилка & Модель пояснює позакласичні ефекти полів первинної зорової кори; окремі спостереження узгоджуються (огляд), як загальну будову конвеєра не доведено & \cite{RaoBallard1999,Keller2018} & гіпотеза \\
10 & Складні картинки потребують зворотних зв'язків і кількох кіл & Мавпа: для «складних» картинок рішення в нижньоскроневій корі складається на $\sim30$\,мс пізніше; ці пізні відповіді погано передбачає пряма мережа й краще --- мережі зі зворотними зв'язками або глибші & \cite{Kar2019} & виміряно \\
11 & Непомітна підробка пікселів обманює мережу; зір людини набагато стійкіший, але не невразливий & На мережах: мале спеціально дібране збурення змінює відповідь; приклад «панда $\to$ гібон» --- із другої праці. У людини за короткого показу підробки, що добре переносяться між мережами, зсувають вибір & \cite{Szegedy2014,Goodfellow2015,Elsayed2018} & виміряно \\
12 & Ілюзії очікування: ненадійний вхід поступається очікуванню --- бліде рухається «повільніше», сукню бачать по-різному (підрозділ~\ref{sec:illusions}) & Ґратки меншого контрасту здаються рухомими повільніше. Опитування $1401$ людини: сукню бачать у трьох стійких варіантах, підказка про освітлення перемикає колір; у $\sim13\,000$ спостерігачів вирішує припущення про освітлення. Баєсова модель з очікуванням повільних рухів пояснює низку ілюзій руху & \cite{Thompson1982,LaferSousa2015,Wallisch2017,Weiss2002,Purves2011} & гіпотеза \\
13 & Регулювання підсилення: водоспад, нахил і контраст; ґратка Германа --- не гальмування сусідів & Детектори напрямку сітківки кроля за тривалого руху знижують частоту, а після зупинки падають нижче за фон. Ілюзію нахилу відтворює модель із діленням на активність сусідів. Зміни ґратки, що не змінюють відповіді вихідних нейронів сітківки, прибирають плями --- пояснення сусідським гальмуванням відкинуто & \cite{BarlowHill1963,Schwartz2009,SchillerCarvey2005} & виміряно \\
14 & Компенсація затримок: рухомий предмет здається попереду спалаху & Ілюзію виміряно. Психофізика суперечить і продовженню руху вперед, і різниці затримок; запропоновано пояснення заднім числом --- за подіями $\sim80$\,мс після спалаху. Суперечку не розв'язано & \cite{Nijhawan1994,EaglemanSejnowski2000} & гіпотеза \\
15 & Нерухомі «змії» обертаються: різниця затримок~\eqref{eq:latency} читається детекторами напрямку & Ілюзія працює й без кольору; пари сусідніх елементів породжують її поодинці; нейрони кори мавпи, вибіркові до напрямку, дають спрямовану відповідь на ці нерухомі пари. У людини обертання запускають мікрострибки очей і кліпання & \cite{Conway2005,OteroMillan2012} & виміряно \\
16 & Двозначні картинки: взаємне гальмування, втома й шум; зміни викликає головно шум & Модель конкурентних груп нейронів: зміни в ній викликає \emph{шум}, без шуму їх немає; вона пояснює зростання частоти змін із силою стимулу. Втома тут відіграє меншу роль & \cite{MorenoBote2007} & гіпотеза \\
17 & Частина ілюзій --- наслідок задачі, якої вчиться машина & Мережа, навчена передбачати наступний кадр відео, «бачить» обертання нерухомих кілець і не бачить його на контрольних картинках; мережі для очищення й підвищення різкості зображень піддаються ілюзіям кольору та яскравості, як людина & \cite{Watanabe2018,GomezVilla2020} & виміряно \\
\bottomrule
\end{longtable}}

\section{Розділ~\ref{sec:muscles}. Вихід на м'язи}

Будову виходу --- рухові одиниці, запас надійності синапсу, увімкнення за розміром, пару м'язів, петлю розтягу й генератори ритму --- виміряно прямо; умова стійкості петлі із затримкою --- теорема; внутрішня модель тіла --- добре обґрунтована гіпотеза; числа рисунків --- розрахунок за моделлю книги.

{\footnotesize
\setlength{\tabcolsep}{3pt}\renewcommand{\arraystretch}{1.15}
\begin{longtable}{>{\raggedright}p{0.5cm}>{\raggedright}p{4.0cm}>{\raggedright}p{5.6cm}>{\raggedright}p{2.3cm}>{\raggedright\arraybackslash}p{1.7cm}}
\toprule
№ & Твердження & Дослід і що виміряно & Джерело & Статус \\
\midrule\endfirsthead
\toprule
№ & Твердження & Дослід і що виміряно & Джерело & Статус \\
\midrule\endhead
1 & Рухова одиниця: один \nt{ACET}-нейрон керує групою волокон, від кількох сотень до пари тисяч; у м'язах тонкого підлаштування одиниці дрібніші & М'язи людини, підрахунок рухових аксонів і м'язових волокон: у середньому близько $340$ волокон на нейрон у міжкістковому м'язі кисті, близько $410$ у плечопроменевому м'язі, близько $1930$ у медіальному литковому. Точного числа для найдрібніших одиниць розділ не наводить. & \cite{Feinstein1955mu} & виміряно \\
2 & Синапс на м'язі вивільняє в кілька разів більше медіатора, ніж потрібно для спрацювання волокна & Огляд вимірювань на нервово-м'язових синапсах: запас надійності (відношення вивільненого медіатора до порогового) у дорослих ссавців зазвичай $3$--$5$; у людини запас тримається більше на складках мембрани одержувача, ніж на величині вивільнення. & \cite{WoodSlater2001} & виміряно \\
3 & Посмикування~\eqref{eq:twitch} наростає за час $T_c$ порядку $50\ms$ & Поодинокі рухові одиниці кисті людини за довільного зусилля, силу усереднено за імпульсами одиниці: час наростання $30$--$100\ms$, у $80\,\%$ одиниць менше за $70\ms$. Функція $(t/T_c)e^{1-t/T_c}$ --- апроксимація з моделі пулу одиниць. & \cite{MilnerBrown1973rec, Fuglevand1993} & виміряно \\
4 & Сума посмикувань~\eqref{eq:forcesum}: за $5$, $15$ і $40$ імп/с сила $0.2$--$1.1F_1$, $1.8$--$2.2F_1$ і близько $5.4F_1$ (рис.~\ref{fig:tetanus}) & Розрахунок за \texttt{models/\allowbreak muscle\_\allowbreak models.py}, $T_c=50\ms$: $0.21$--$1.10$, $1.76$--$2.19$ і $5.28$--$5.49\,F_1$ на останніх $0.3\,\text{с}$. Лінійне додавання --- спрощення: насичення справжнього м'яза в моделі немає. & \texttt{models/\allowbreak muscle\_\allowbreak models.py} & модель \\
5 & Одиниці вмикаються за розміром; найслабші в сотню разів слабші за найсильніші & Вентральні корінці кішки: за наростаючого рефлекторного входу першими розряджаються нейрони з дрібними імпульсами в корінці, тобто малі. Міжкістковий м'яз кисті людини: сила посмикування одиниць від $0.1$ до $10$~г і зростає майже лінійно із силою, за якої одиниця вмикається. & \cite{Henneman1957, MilnerBrown1973rec} & виміряно \\
6 & Крок сили пропорційний силі, $\Delta F/F\approx q-1$~\eqref{eq:sizeprinciple}: рухома кома & Виведено в розділі з геометричної прогресії сил. Узгоджується з дослідом: за великих зусиль на той самий приріст сили вмикається набагато менше нових одиниць. & виведення в розділі; \cite{MilnerBrown1973rec} & теорія \\
7 & Розкид сили зростає пропорційно самій силі & Довільне ізометричне зусилля м'яза великого пальця: стандартне відхилення сили лінійно зростає із середньою силою; за того самого зусилля, викликаного електричною стимуляцією, розкид сталий. Модель пулу: лінійне зростання дає ввімкнення одиниць за порядком сили. & \cite{Jones2002sdn} & виміряно \\
8 & Плавність рухів --- наслідок шуму, що залежить від сигналу & Модель: траєкторія, що мінімізує розкид кінцевої точки за такого шуму, відтворює виміряні траєкторії ока й руки. Прямого досліду, який відрізняв би цю причину від інших, немає. & \cite{HarrisWolpert1998} & гіпотеза \\
9 & Різниця сил пари м'язів задає момент, сума --- жорсткість~\eqref{eq:antagonists} & Теорія керування імпедансом через спільне напруження. Рука людини: жорсткість кожного суглоба, виміряна малими поштовхами, приблизно пропорційна моменту в ньому; різні співвідношення спільного напруження змінюють форму й напрямок еліпса жорсткості кисті. & \cite{Hogan1984, GomiOsu1998stiff} & виміряно \\
10 & Датчик розтягу збуджує свій \nt{ACET}-нейрон і через гальмівного посередника вимикає антагоніста (рис.~\ref{fig:antagonists}) & Спинний мозок кішки: поодинокий посередник, що збуджується волокнами датчиків розтягу, дає моносинаптичні гальмівні потенціали в \nt{ACET}-нейронах м'яза-антагоніста, синаптична затримка $0.28$--$0.42\ms$. Медіатор посередника (гліцин) у цьому досліді не визначали. & \cite{Jankowska1972Ia} & виміряно \\
11 & Затримка петлі: спинномозкової $25$--$30\ms$, через кору в півтора--три рази більша, через око $100$--$200\ms$ & Рука людини, поштовх по суглобу: коротка відповідь м'язів через $20$--$45\ms$, довга $45$--$105\ms$, довільна $120$--$180\ms$. Зміщення видимого положення пальця під час руху: поправка через $160\ms$ у середньому. & \cite{Pruszynski2008rapid, SaundersKnill2003vis} & виміряно \\
12 & $\dd x/\dd t=-Gx(t-d)$ стійке при $Gd<\pi/2$~\eqref{eq:delaystable}; на межі частота $1/(4d)$ & Теорема про корені рівняння із запізненням. Перевірка за \texttt{models/\allowbreak muscle\_\allowbreak models.py}, $d=30\ms$: до $3\,\text{с}$ амплітуда $3\cdot10^{-6}$ при $G=40$, $0.13$ при $G=50$, $38$ при $G=55$ за секунду (межа $52$). & \cite{Hayes1950} & теорія \\
13 & Фізіологічне тремтіння $8$--$12$~Гц; петля розтягу --- одне з його джерел & Огляд вимірювань: дрібне тремтіння в діапазоні близько $10$~Гц є у здорових людей; його походження змішане --- центральні коливання, властивості розрядів одиниць, механічний резонанс і резонанс рефлекторної петлі, і за різних умов переважає різне. & \cite{McAuleyMarsden2000} & гіпотеза \\
14 & Мозок передбачає результат команди за внутрішньою моделлю тіла & Людина тримає вантаж щипком і рухає рукою: за інерційного, в'язкого й пружного навантаження сила стискання змінюється одночасно з навантаженням, а не із запізненням петлі, тобто навантаження передбачено. Огляд зводить такі досліди; прямого спостереження самої моделі в нейронах немає. & \cite{FlanaganWing1997grip, WolpertGhahramani2000} & гіпотеза \\
15 & Ритм ходьби, дихання, жування породжують місцеві мережі за сталого входу & Огляд дослідів: ізольована нервова система (спинний мозок, ганглії) видає ритмічний руховий візерунок без ритмічного входу й без зворотного зв'язку від м'язів. & \cite{MarderBucher2001} & виміряно \\
16 & Взаємне гальмування з утомою~\eqref{eq:halfcentre} дає ритм із періодом $0.84\,\text{с}\approx2\tau_a$ (рис.~\ref{fig:halfcentre}) & Теорема: мережа із взаємним гальмуванням і адаптацією без стійкого стану за сталого входу обов'язково коливається. Розрахунок за \texttt{models/\allowbreak muscle\_\allowbreak models.py} ($I=1$, $\beta=2$, $b=2$, $\tau=20\ms$, $\tau_a=0.4\,\text{с}$): період $0.84\,\text{с}$. Що справжні генератори влаштовані саме так, не показано. & \cite{Matsuoka1985osc} & модель \\
\bottomrule
\end{longtable}}

\section{Розділ~\ref{sec:pain}. Біль}

Датчики болю, два проводи, ворота в спинному мозку, низхідну ручку гучності, сенсибілізацію й захоплення уваги виміряно прямо; формули воріт і сенсибілізації --- спрощені моделі книги; правило, за яким машина вибирає гучність тривоги, --- гіпотеза.

{\footnotesize
\setlength{\tabcolsep}{3pt}\renewcommand{\arraystretch}{1.15}
\begin{longtable}{>{\raggedright}p{0.5cm}>{\raggedright}p{4.0cm}>{\raggedright}p{5.6cm}>{\raggedright}p{2.3cm}>{\raggedright\arraybackslash}p{1.7cm}}
\toprule
№ & Твердження & Дослід і що виміряно & Джерело & Статус \\
\midrule\endfirsthead
\toprule
№ & Твердження & Дослід і що виміряно & Джерело & Статус \\
\midrule\endhead
1 & Високий поріг: пороги нагрівання в різних датчиків болю від $41^\circ$ до $46$--$48^\circ$ & Запис поодиноких волокон. Шкіра мавпи: медіанний поріг повільних (C) датчиків $41^\circ$, швидких датчиків другого типу $46^\circ$, першого типу --- вище за $53^\circ$. Шкіра людини: C-датчики, що відповідають і на тиск, $40.7^\circ$, що не відповідають на тиск --- $48^\circ$. & \cite{Treede1995heat, Weidner1999cfib} & виміряно \\
2 & Датчики болю звикають набагато слабше за датчики дотику й після легкого ушкодження стають чутливішими; послаблення після сильного нагрівання виміряно в одного типу & Легкий опік шкіри ($50^\circ$, $100\,\text{с}$): через $5$--$10$~хв поріг болю в людей і поріг C-датчиків у мавпи знижено на $2$--$6^\circ$; в інших шкірних датчиків такого зсуву немає. У швидких датчиків другого типу відповідь після сильного нагрівання пригнічено. Порівняння звикання з датчиками дотику --- загальна оцінка, окремим дослідом тут не підкріплена. & \cite{LaMotte1982hyper, Treede1995heat} & виміряно \\
3 & Два проводи: швидкий $5$--$30$~м/с, повільний $0.5$--$2$~м/с; час приходу $t=L/v$~\eqref{eq:paintime} & Швидкість проведення: швидкі датчики болю мавпи в середньому $15$ і $25$~м/с (два типи); C-датчики людини $0.8$--$1.0$~м/с. При $L=1$~м це десятки мілісекунд і близько секунди. & \cite{Treede1995heat, Weidner1999cfib} & виміряно \\
4 & Біль приходить двічі: спершу швидкий і точний, потім тупий і довгий & Час реакції на нагрівання передпліччя людини: від $1100$ до $400\ms$ зі зростанням температури; сильне нагрівання волосистої шкіри дає подвійний біль, долоні --- ні. Перший біль можуть нести лише волокна, швидші за $6$~м/с, --- за латентністю йому відповідають швидкі датчики другого типу, яких у долоні немає. Розподіл ролей («де» і «наскільки погано») --- тлумачення. & \cite{CampbellLaMotte1983first, Treede1995heat} & виміряно \\
5 & Ворота: вихідний нейрон спинного мозку отримує біль і дотик, гальмівний посередник дотиком збуджується (рис.~\ref{fig:gate}) & Схему воріт запропоновано як теорію. Миші, генетична мітка й видалення груп нейронів: збуджувальні нейрони із соматостатином потрібні для механічного болю; гальмівні нейрони з динорфіном не дають входу від датчиків дотику дійти до них --- без цих нейронів легкий дотик викликає біль. & \cite{MelzackWall1965, Duan2014} & виміряно \\
6 & Дотик приглушує біль & Людина, лазерні імпульси болю на руці: одночасний дотик знижує оцінку болю й здатність розрізняти енергію імпульсів; ефект слабшає з відстанню між точкою дотику й точкою болю в межах одного шкірного сегмента. & \cite{Mancini2014touch} & виміряно \\
7 & Припущення теорії воріт: сильний біль сам вимикає гальмівного посередника, і ворота розчахуються & У розділі позначено як припущення вихідної теорії. Праця на мишах описує, як ворота не пускають дотик у канал болю; вимкнення посередника болем у ній не показано. & \cite{MelzackWall1965} & гіпотеза \\
8 & Ворота~\eqref{eq:gate}: поріг $0.18$ без дотику, $0.86$ з дотиком, $0.11$ при $g=2$; при $\nu=1$ дотик знижує вихід з $1$ до $0.25$, при $\nu=2$ не діє; дотик сам не проходить (рис.~\ref{fig:gatecurves}) & Розрахунок за \texttt{models/\allowbreak pain\_\allowbreak gate.py} з вагами з тексту відтворює всі ці числа. & \texttt{models/\allowbreak pain\_\allowbreak gate.py} & модель \\
9 & Низхідні шляхи зі стовбура змінюють підсилення воріт в обидва боки & Щур, довгастий мозок, запис під час нагрівання хвоста: одні нейрони розряджаються перед відсмикуванням («увімк»), інші замовкають («вимк»); слабка стимуляція цієї ділянки пригнічує рефлекс. Огляд: ті самі кола і полегшують, і гальмують передачу болю, опіоїди діють через них. & \cite{Fields1983rvm, Fields2004} & виміряно \\
10 & Очікування полегшення приглушує біль через опіоїдний канал & Люди з гострим болем: у тих, кому очікування полегшення зменшило біль, блокада опіоїдних рецепторів повернула біль до рівня решти; в інших блокада болю не змінювала. & \cite{Levine1978} & виміряно \\
11 & Гучність тривоги мозок вибирає за вигодою переривання & Досліду, в якому виміряно правило вибору гучності, немає. & --- & гіпотеза \\
12 & Сенсибілізація: після ушкодження вихід воріт зростає, поріг падає, почасти за рахунок спинного мозку; вона тримається й після припинення ушкоджувального стимулу & Щур: після ушкодження шкіри поріг згинального рефлексу падає, а відповідь зростає; електрофізіологічний аналіз показує, що частина зростання виникає в самому спинному мозку. Ушкоджувальний стимул викликає тривалі порушення чутливості, що переживають сам стимул. Точного часу спаду розділ не наводить. & \cite{Woolf1983} & виміряно \\
13 & Сенсибілізація~\eqref{eq:sensitization} --- додатний зворотний зв'язок; за сильної петлі тривога може защепнутися після загоєння & Випливає з моделі~\eqref{eq:sensitization} (задача наприкінці розділу). Досліду, який показав би, що стійкий біль після загоєння --- защепнута петля саме такого вигляду, немає. & виведення в розділі & гіпотеза \\
14 & Біль --- від'ємна винагорода, і навчання на ньому йде за помилкою часових різниць & Томограф, люди, ланцюжки провісників болю: активність вентрального стріатума й передньої острівцевої кори збігається із сигналами навчання за часовими різницями. Мавпа: частина \nt{DOPA}-нейронів збуджується від неприємного обдування та його провісника, інша частина гальмується. & \cite{Seymour2004, Matsumoto2009} & виміряно \\
15 & Біль перериває поточну роботу & Люди, електричний больовий стимул і звукові проби: відповідь на пробу через $100\ms$ після початку болю помітно сповільнена, через $1.5\,\text{с}$ різниці з контролем уже немає. Огляд зводить такі досліди в модель захоплення уваги. & \cite{Crombez1996disrupt, EcclestonCrombez1999} & виміряно \\
16 & Відсмикування замикається в спинному мозку & Щур: нейрони глибоких шарів заднього рогу повторюють просторовий візерунок рефлексу відсмикування окремих м'язів; їх не вдається збудити антидромно з шийного відділу, тобто це спінальні посередники. & \cite{Schouenborg1995wr} & виміряно \\
\bottomrule
\end{longtable}}

\section{Розділ~\ref{sec:motorlearning}. Навчання рухів}

Правило найменших квадратів і вигода окремого проводу помилки --- теореми; будову схеми з проводом помилки, послаблення за збігу, підлаштування сліду під затримку, післядію й спонтанне повернення навички виміряно; те, що вся схема працює як навчання з учителем і будує внутрішню модель, --- провідна гіпотеза; числа двопроцесної моделі --- розрахунок за моделлю книги.

{\footnotesize
\setlength{\tabcolsep}{3pt}\renewcommand{\arraystretch}{1.15}
\begin{longtable}{>{\raggedright}p{0.5cm}>{\raggedright}p{4.0cm}>{\raggedright}p{5.6cm}>{\raggedright}p{2.3cm}>{\raggedright\arraybackslash}p{1.7cm}}
\toprule
№ & Твердження & Дослід і що виміряно & Джерело & Статус \\
\midrule\endfirsthead
\toprule
№ & Твердження & Дослід і що виміряно & Джерело & Статус \\
\midrule\endhead
1 & Правило~\eqref{eq:lms} у середньому йде за градієнтом квадрата помилки & Результат теорії адаптивних фільтрів: поправка, пропорційна входу й помилці виходу, є стохастичним спуском за градієнтом середнього квадрата помилки. & \cite{WidrowHoff1960} & теорія \\
2 & З проводом помилки на кожен вихід швидкість не залежить від кількості виходів; з одним спільним числом --- падає з кількістю нейронів, що шумлять (твердження~\ref{prop:teachers}) & Криві навчання лінійної мережі: навчання за випадковими збуреннями нейронів повільніше за точний градієнт у стільки разів, скільки нейронів збурюють. & \cite{Werfel2005} & теорія \\
3 & Схема з проводом помилки працює як навчання з учителем (твердження~\ref{prop:errorwire}) & Теорії, виведені з будови схеми. Досліди рядків 7--10 з ними узгоджуються; що вся схема працює саме так, не доведено. & \cite{Marr1969, Albus1971} & гіпотеза \\
4 & Розширювач: близько $7\cdot10^{10}$ дрібних \nt{GLUT}-нейронів, більше, ніж у всьому іншому мозку & Підрахунок ядер у суспензії розчиненої тканини: у мозочку людини $69\cdot10^9$ нейронів із $86\cdot10^9$ в усьому мозку. Стереологія: $101\cdot10^9$ дрібних нейронів у мозочку літніх чоловіків. & \cite{Azevedo2009, Andersen1992cereb} & виміряно \\
5 & Підняття розмірності входу робить потрібну залежність лінійною & Теорема про кількість розбиттів: зважена сума $n$ входів із порогом реалізує довільну розмітку приблизно $2n$ векторів загального положення. Що мозочок користується саме цим, --- частина гіпотези рядка~3. & \cite{Cover1965} & теорія \\
6 & Близько $3\cdot10^7$ вихідних \nt{GABA}-нейронів, у кожного порядку $10^5$ входів & Стереологія мозочка людини: $30.5\cdot10^6$ вихідних нейронів. Електронна мікроскопія, щур: близько $1.75\cdot10^5$ входів від розширювача на один вихідний нейрон. Гальмівний медіатор вихідних нейронів --- в огляді. & \cite{Andersen1992cereb, NapperHarvey1988, Ito2001} & виміряно \\
7 & Рівно один провід помилки на вихідний нейрон, близько разу на секунду, щоразу потужний спалах & Огляд записів: кожен вихідний нейрон отримує один ліаноподібний \nt{GLUT}-вхід, що викликає складний імпульс із частотою порядку $1$~Гц. & \cite{Ito2001} & виміряно \\
8 & Імовірність спалаху залежить від напрямку помилки руху & Мавпа, окорухова частина мозочка, адаптація саккад: напрямок зорової помилки задає ймовірність складного імпульсу, величина --- його час; спалах змінює прості імпульси в наступній спробі й зсуває рух уздовж переважної помилки клітини. & \cite{Herzfeld2018} & виміряно \\
9 & Збіг спалаху помилки з активним входом послаблює вхід ($-\eta e_i x_j$) & Огляд дослідів: спільна стимуляція входу від розширювача й проводу помилки викликає довготривале послаблення цього входу; стимуляція одного входу --- ні. & \cite{Ito2001} & виміряно \\
10 & Довжину сліду підлаштовано під затримку помилки: за затримки $100\ms$ найсильніше послаблюється вхід, що був активним за $100\ms$ & Зрізи й живі миші: у частині, що відповідає за окорухове навчання, послаблення вузько налаштоване на інтервал близько $120\ms$ між входом і помилкою --- стільки йде сигнал помилки в цій задачі; в іншій частині різні клітини налаштовано на різні інтервали. & \cite{Suvrathan2016} & виміряно \\
11 & Схема --- адаптивний фільтр~\eqref{eq:adaptivefilter}, що вивчає внутрішню модель тіла & Модель, виведена з будови схеми. Прямого досліду, де вивчений фільтр виміряно як передавальну функцію тіла, немає. & \cite{Fujita1982} & гіпотеза \\
12 & Передбачення віднімає наслідки власних рухів, тому себе не полоскочеш & Люди: власний дотик здається менш лоскітним, ніж такий самий чужий; у томографі соматосенсорна кора відповідає на власний дотик слабше, а мозочок менш активний, коли рух торкається шкіри. Пояснення через віднімання передбачення --- гіпотеза. & \cite{Blakemore1998} & гіпотеза \\
13 & За зовнішньої сили промах зменшується, а після її зняття рука промахується в інший бік & Люди тягнуться рукояттю робота в силовому полі: траєкторії повертаються до прямих; після зняття поля післядія --- майже дзеркальне відображення перших викривлень; вона переноситься й у нетреновані частини простору. & \cite{ShadmehrMussaIvaldi1994} & виміряно \\
14 & Навчаються два процеси~\eqref{eq:tworate}; після зворотного збурення навичка повертається сама & Люди, силове поле, далі коротке зворотне поле й спроби із затиснутою помилкою: поправка сама повертається до старої; модель із двох процесів це відтворює, модель з одного --- ні. & \cite{Smith2006} & виміряно \\
15 & Числа рис.~\ref{fig:tworate}: $0.84$ (повільний $0.63$) за $200$ рухів; $6$ зворотних; швидкий $-0.53$, повільний $0.46$; повернення до $0.37$ за $40$ рухів & Розрахунок за \texttt{models/\allowbreak motor\_\allowbreak learning.py}, $A_f=0.92$, $B_f=0.1$, $A_s=0.996$, $B_s=0.02$: $0.840$ ($0.634$ і $0.205$), $6$ рухів, $-0.531$ і $0.457$, пік $0.371$ на $40$-му русі. & \texttt{models/\allowbreak motor\_\allowbreak learning.py} & модель \\
16 & Два процеси потрібні, бо тіло змінюється з різними швидкостями & Теорія: оптимальна оцінка за завад із кількома часами життя дає кілька фільтрів із різними сталими часу й пояснює спонтанне повернення та прискорене повторне навчання. & \cite{Kording2007} & гіпотеза \\
\bottomrule
\end{longtable}}

\section{Розділ~\ref{sec:chemoutput}. Хімічний вихід}

Два проводи до серця та їхню різну швидкість, від'ємний зворотний зв'язок у гормональних каскадах, код частотою порцій, добовий генератор і його підстроювання світлом виміряно прямо; межа $K<8$ --- теорема теорії автоматичного керування, виведена в розділі; пульсації, породжувані самим зворотним зв'язком каскаду, --- гіпотеза із сильним дослідом на її користь.

{\footnotesize
\setlength{\tabcolsep}{3pt}\renewcommand{\arraystretch}{1.15}
\begin{longtable}{>{\raggedright}p{0.5cm}>{\raggedright}p{4.0cm}>{\raggedright}p{5.6cm}>{\raggedright}p{2.3cm}>{\raggedright\arraybackslash}p{1.7cm}}
\toprule
№ & Твердження & Дослід і що виміряно & Джерело & Статус \\
\midrule\endfirsthead
\toprule
№ & Твердження & Дослід і що виміряно & Джерело & Статус \\
\midrule\endhead
1 & Ритмоводій серця сам б'ється близько $100$ разів на хвилину; у спокої гальмо натиснуте, і частота зазвичай порядку $60$--$70$ & Люди, повне вимкнення обох проводів до серця: власна частота вища за частоту спокою і спадає з віком, у дорослих порядку $100$ ударів на хвилину. Отже, у спокої переважає гальмо. Частота спокою $60$--$70$ --- типове значення, окремо не цитоване. & \cite{JoseCollison1970ihr} & виміряно \\
2 & Газ \nt{NORA} і гальмо \nt{ACET} --- фільтри нижніх частот, гальмо швидше~\eqref{eq:gasbrake} & Собака, частотно-модульована стимуляція блукаючого або симпатичного нерва серця й передатна функція до частоти передсердь: обидва шляхи --- фільтри нижніх частот, у симпатичного частота зрізу значно нижча й додано чисту затримку близько $1.7\,\text{с}$. Характеристики залежать від середнього тонусу, тож лінійна формула~\eqref{eq:gasbrake} --- наближення. & \cite{Berger1989} & виміряно \\
3 & Гальмо швидке, бо калієвий канал \nt{ACET} відкриває без другого посередника, а \nt{NORA} --- через ланцюжок реакцій & Мембранні клаптики клітин передсердь: калієвий канал, що відкривається ацетилхоліном, відкривають $\beta\gamma$-субодиниці G-білка, пов'язаного з рецептором, без другого посередника всередині клітини. Шлях \nt{NORA} через цАМФ тут не цитовано. & \cite{Logothetis1987gbg} & виміряно \\
4 & Кров обходить тіло за час порядку хвилини; \nt{ADRE} з кровотоку доходить до всіх органів із рецептором & Загальна оцінка порядку величини; у розділі так і сформульовано, джерело не цитовано. & --- & оцінка \\
5 & Гормональний каскад охоплено від'ємним зворотним зв'язком: кінцевий гормон гальмує перші ступені & Огляд дослідів: гормони кори надниркових залоз пригнічують викид АКТГ гіпофізом і рилізинг-гормону гіпоталамусом у трьох часових діапазонах --- швидкому, проміжному й повільному. & \cite{KellerWood1984fb} & виміряно \\
6 & Три однакові ступені стійкі за $K<8$~\eqref{eq:cascadestable}, на межі період $2\pi\tau/\sqrt3\approx3.6\tau$ & Виведено в розділі: на частоті $\sqrt3/\tau$ три ступені дають зсув фази $180^\circ$ і ослаблення у $8$ разів. & виведення в розділі & теорія \\
7 & Похибка рівня $1/(1+K)$, тому точніше приблизно за $10\,\%$ такий регулятор не тримає (твердження~\ref{prop:cascade}) & Наслідок рядка~6 для пропорційного зворотного зв'язку. Справжня вісь має зворотний зв'язок на кількох ступенях і в кількох часових діапазонах (рядок~5), тож межа стосується моделі~\eqref{eq:cascade}, а не виміряна. & виведення в розділі & теорія \\
8 & За $K=4$ і $7$ рівень усталюється, за $K=12$ --- рівні пульсації з періодом $3.7$--$3.9\,\tau$ (рис.~\ref{fig:cascade}) & Розрахунок за \texttt{models/\allowbreak chem\_\allowbreak models.py}: за $K=12$ послідовні періоди $3.9$, $3.7$, $3.8$, $3.7\,\tau$; за $K=4$ розмах наприкінці рахунку $0.003$. & \texttt{models/\allowbreak chem\_\allowbreak models.py} & модель \\
9 & Гормон стресу виходить пульсами приблизно раз на годину; пульси породжує сам зворотний зв'язок, без окремого генератора & Щури, постійне вливання рилізинг-гормону: АКТГ і кортикостерон коливаються з фізіологічною, майже годинною частотою --- ритмічний вхід від гіпоталамуса для пульсів не потрібен. Модель гіпофіз--наднирник із запізнілим зворотним зв'язком дає такі коливання. & \cite{Walker2012glc, Walker2010} & гіпотеза \\
10 & Отримувач читає частоту порцій, а не рівень & Мавпи без власного викиду рилізинг-гормону гонадотропінів: постійне вливання не відновлює викиду гормонів гіпофіза, порції раз на годину відновлюють; перехід на постійне вливання знову глушить відповідь. Утома рецептора як фільтр верхніх частот --- тлумачення. & \cite{Belchetz1978} & виміряно \\
11 & Різні частоти порцій по-різному діють на різні клітини однієї залози & Ті самі мавпи: почастішання порцій до $2$--$5$ на годину знижує обидва гормони гіпофіза; порідшання до однієї за $3$~години знижує один (ЛГ) і підвищує другий (ФСГ), тож їхнє відношення задається частотою. & \cite{Wildt1981gnrh} & виміряно \\
12 & Добовий ритм породжує внутрішньоклітинний від'ємний зворотний зв'язок із затримкою в кілька годин & Огляд: білки годинникових генів із затримкою пригнічують транскрипцію власних генів; петля з транскрипції й трансляції дає період близько доби. & \cite{TakahashiJS2017} & виміряно \\
13 & Головний генератор --- група з десятків тисяч нейронів, що синхронізує решту тіла & Огляд: супрахіазматичне ядро --- головний добовий генератор; окремі його нейрони в культурі --- самостійні генератори, мережа їх синхронізує; у миші близько $10^4$ нейронів із кожного боку. & \cite{Welsh2010scn} & виміряно \\
14 & Власний період у людини $24.18$~години & Люди в тьмяному світлі за розкладом, відірваним від доби: період ритмів мелатоніну, температури й кортизолу в середньому $24.18$~години в молодих і літніх, із вузьким розкидом. & \cite{Czeisler1999} & виміряно \\
15 & Світло підводить генератор як фазове автопідстроювання~\eqref{eq:pll}; потрібен зсув близько $11$~хв на добу & Люди, одиночне яскраве світло $6.7$~години в різний час доби: крива фазової відповіді першого типу з розмахом $5$~годин --- затримка до мінімуму температури тіла, випередження після. Рівняння~\eqref{eq:pll} --- стандартна модель; $11$~хв $=0.18$~години --- арифметика. & \cite{Khalsa2003prc} & виміряно \\
16 & Без світла підстроювання немає, і ритм переходить на власний період & Цілком незрячі люди без сприйняття світла, що живуть у звичайному суспільстві: приблизно в половини ритми мелатоніну й кортизолу йдуть зі своїм періодом, що не збігається з добою. & \cite{Sack1992blind} & виміряно \\
\bottomrule
\end{longtable}}

\section{Розділ~\ref{sec:debugging}. Налагодження машини}

Що чує електрод, маловимірність активності, роботу інтерфейсів на жорстких ґратках, спільне навчання мозку й декодера та зорові точки від стимуляції виміряно в рецензованих дослідах; оцінки потоків даних --- арифметика розділу; про вживлений людям пристрій Neuralink, наскільки вдалося перевірити, публікувала дані переважно сама компанія.

{\footnotesize
\setlength{\tabcolsep}{3pt}\renewcommand{\arraystretch}{1.15}
\begin{longtable}{>{\raggedright}p{0.5cm}>{\raggedright}p{4.0cm}>{\raggedright}p{5.6cm}>{\raggedright}p{2.3cm}>{\raggedright\arraybackslash}p{1.7cm}}
\toprule
№ & Твердження & Дослід і що виміряно & Джерело & Статус \\
\midrule\endfirsthead
\toprule
№ & Твердження & Дослід і що виміряно & Джерело & Статус \\
\midrule\endhead
1 & Електрод чує імпульси нейронів у радіусі порядку сотні мікрометрів, зазвичай від одного до кількох; їх розрізняють за формою & Щур, поле CA1, одночасний внутрішньоклітинний і позаклітинний запис: форма позаклітинного імпульсу повторює ширину й амплітуду внутрішньоклітинного. Оцінка: тетрод міг би розрізняти $60$--$100$ нейронів, а насправді виділяють близько шести. & \cite{Henze2000extra} & виміряно \\
2 & У мозку $8.6\cdot10^{10}$ нейронів; щуп чує приблизно одну стомільйонну & Підрахунок ядер у суспензії розчиненої тканини: $86\cdot10^9$ нейронів у мозку людини. Частка --- арифметика розділу. & \cite{Azevedo2009} & виміряно \\
3 & Активність сотень нейронів рухової кори вкладається в один-два десятки спільних змінних & Огляд записів у руховій корі мавп: активність популяції описується небагатьма прихованими змінними, спільними для багатьох нейронів. & \cite{Gallego2017} & виміряно \\
4 & Шум сусідів узгоджений, і сотня нейронів несе майже стільки ж, скільки тисяча (твердження~\ref{prop:pooling}) & Зорова кора мавпи: коефіцієнт кореляції шуму сусідніх нейронів близько $0.1$, і виграш від усереднення насичується на сотні нейронів. Теорія кореляцій, що обмежують інформацію. & \cite{Zohary1994, MorenoBote2014} & виміряно \\
5 & Сирий потік $2\cdot10^8$~біт/с проти $8\cdot10^4$~біт/с в імпульсах~\eqref{eq:probedata} & Арифметика розділу за ємністю потоку імпульсів~\eqref{eq:spikeentropy}. Параметри оцифрування реальні: кристал Neuralink --- $19.3$~кГц, $10$~біт на канал. & виведення в розділі; \cite{Rieke1997, Musk2019} & теорія \\
6 & Перша система Neuralink: кристали підсилюють і оцифровують, сирий потік іде кабелем, імпульси виділяє програма ззовні; виділення на кристалі, закритий корпус, радіо й бездротове живлення --- у пристрою для людей, за повідомленнями компанії & Стаття компанії: повний сирий потік кабелем USB-C, імпульси виділяє програма в реальному часі поза кристалом; стиснення на борту, бездротові живлення й передавання названо планом для клінічних пристроїв. Для вживленого людям пристрою --- лише повідомлення компанії. Що виділяти імпульси на кристалі необхідно, --- висновок розділу з~\eqref{eq:probedata}. & \cite{Musk2019}; звіти Neuralink, рецензованої статті немає & виміряно (стаття компанії); для людей --- звіт компанії \\
7 & До $3072$ електродів на $96$ нитках по $32$; нитки вводить робот, оминаючи судини & Стаття компанії: $3072$ електроди на $96$ нитках по $32$; робот уводить $6$ ниток ($192$ електроди) за хвилину; до $70\,\%$ електродів у щурів із довготривало вживленою ґраткою чують імпульси. & \cite{Musk2019} & виміряно (стаття компанії) \\
8 & У людей близько тисячі каналів на $64$ нитках; частина ниток змістилася; курсор порядку $10$~біт/с & Лише повідомлення компанії. Пошук у PubMed не знайшов рецензованої статті з результатами на людях; це узгоджується із застереженням у тексті розділу. & звіти Neuralink; рецензованої статті немає & виміряно (звіт компанії) \\
9 & Інтерфейс на ґратці із сотні електродів керує курсором & Людина з паралічем, $96$ електродів у первинній руховій корі: намір рухати рукою змінює імпульси через три роки після травми; декодер дає «нейронний курсор» (пошта, телевізор), керування протезом кисті й роботом-рукою. & \cite{Hochberg2006} & виміряно \\
10 & Письмо «уявною рукою»: близько $90$ знаків на хвилину, менше ніж $8$~біт/с & Людина з паралічем кисті, спроби писати від руки, рекурентний декодер: $90$ знаків на хвилину з точністю $94.1\,\%$ у реальному часі, понад $99\,\%$ після автокорекції. Оцінка в бітах --- арифметика розділу. & \cite{Willett2021} & виміряно \\
11 & Спроби говорити перетворюються на текст зі швидкістю близько $60$ слів на хвилину & Людина, що втратила розбірливе мовлення, ґратки в корі: $62$ слова на хвилину; $9.1\,\%$ помилок у словах за словника з $50$ слів, $23.8\,\%$ за словника зі $125\,000$. & \cite{Willett2023} & виміряно \\
12 & Інтерфейс витягає малу частку ємності: в одного нейрона десятки біт на секунду, у сотні --- тисячі (твердження~\ref{prop:probe}) & Верхня межа~\eqref{eq:spikeentropy} --- теорія; порівняння з рядками~8, 10 --- висновок розділу. Що вузьке місце --- маловимірність і незнання коду, а не провід, прямим дослідом не перевірено. & \cite{Rieke1997} & теорія \\
13 & Мозок і декодер учаться один в одного & Мавпи керують курсором; декодер підлаштовується в замкненій петлі, а нейрони водночас перенавчаються: точність зростає, навичка зберігається й менше страждає від рухів власної руки. Порада розвести швидкості навчання --- інженерне правило, окремого досліду в розділі немає. & \cite{Orsborn2014coad} & виміряно \\
14 & Струм через електрод збуджує навколишні нейрони й проводи без розбору типу & Кора миші й кішки, двофотонний запис кальцію: навіть слабкий струм збуджує рідкісні нейрони, розкидані на відстань до міліметрів, через пряме збудження аксонів в об'ємі діаметром у десятки мікрометрів. Тобто збудження невибіркове, але не суцільне. & \cite{Histed2009stim} & виміряно \\
15 & Стимуляція зорової кори запалює точки; до $16$ точок одночасно дають змогу впізнати деякі літери й межі предметів & Незряча людина, $96$ електродів у зоровій корі на $6$ місяців: поріг одного електрода $66.8\pm36.5$~мкА; одночасна стимуляція груп (до $16$ електродів --- межа стимулятора) знижує поріг і дає розрізненні візерунки, деякі літери й межі предметів упізнаються. & \cite{Fernandez2021} & виміряно \\
16 & Другий напрям Neuralink --- пристрій, що подає сигнали в зорову кору & Лише повідомлення компанії про розроблення; рецензованих даних про його роботу не знайдено. & звіти Neuralink & гіпотеза \\
17 & Концентрації широкомовних медіаторів можна вимірювати хімічним датчиком у тканині & Зрізи мозку щура, швидка циклічна вольтамперометрія на вуглецевому волокні: виміряно викид і зворотне захоплення серотоніну; речовина виходить за межі синапсу. & \cite{Bunin1998} & виміряно \\
\bottomrule
\end{longtable}}

\section{Розділ~\ref{sec:eetypes}. Нейрони як прилади}

Режими роботи окремих клітин виміряно прямо, струмом через електрод і записом напруги; математика інтегратора й поділу на класи --- класична теорія, а те, що мозок користується переналаштуванням режимів для обчислень, лишається гіпотезою.

{\footnotesize
\setlength{\tabcolsep}{3pt}\renewcommand{\arraystretch}{1.15}
\begin{longtable}{>{\raggedright}p{0.5cm}>{\raggedright}p{4.0cm}>{\raggedright}p{5.6cm}>{\raggedright}p{2.3cm}>{\raggedright\arraybackslash}p{1.7cm}}
\toprule
№ & Твердження & Дослід і що виміряно & Джерело & Статус \\
\midrule\endfirsthead
\toprule
№ & Твердження & Дослід і що виміряно & Джерело & Статус \\
\midrule\endhead
1 & Резонатор: повільний струм, що протидіє зміні напруги, робить нейрон смуговим фільтром із максимумом на одиницях--десятках герц (підрозділ~\ref{sec:resonator}) & У зрізах у \nt{GLUT}-нейрони вводили синусоїдальний струм із плавно зростаючою частотою і вимірювали імпеданс $|Z(f)|$: максимум на $2$--$5$~Гц за $33^\circ$C (близько $7$~Гц за $38^\circ$C); резонанс створюють повільні калієвий і катіонний струми, а підсилює постійний натрієвий струм. Огляд дає той самий прийом для багатьох класів клітин & \cite{HuVervaekeStorm2002,HutcheonYarom2000} & виміряно \\
2 & Повільний струм поводиться як індуктивність, разом з ємністю мембрани --- коливальний контур & Підпорогову відповідь аксона на струм виміряно й описано лінеаризованими рівняннями провідностей; еквівалентна схема містить «феноменологічну» індуктивність, величина якої залежить від напруги & \cite{Mauro1970} & теорія \\
3 & Стала часу групи зі зворотним зв'язком $\tau_{\text{еф}}=\tau/(1-w)$~\eqref{eq:perfectint}; за $\tau=0.1$~с і $w=0.995$ --- $20$~с & Виводиться в розділі з лінійного рівняння групи; як механізм утримання положення ока запропоновано в теоретичній роботі & виведення в розділі; \cite{Seung1996} & теорія \\
4 & Інтегратор положення ока тримає вхід мережею, а не властивістю окремої клітини & Внутрішньоклітинний запис у риби у вузлі, що тримає положення ока: після швидкого повороту частота нейрона сходинкою змінюється й тримається; короткий струм в одну клітину спричиняє лише короткочасний зсув, а сходинки напруги зберігаються й нижче порога --- їх підтримують синаптичні входи & \cite{Aksay2001} & виміряно \\
5 & Інтеграторові без витоку потрібне постійне підлаштування навчанням; коли його розладнано, він або забуває, або йде в насичення & Рибу навчали зоровим зворотним зв'язком, пропорційним $\pm$ положенню ока: утримання ставало нестійким або з витоком, стала часу нейронів падала більш ніж на порядок, до $<1$~с; звичайний зоровий зворотний зв'язок поступово повертав стійкість & \cite{Major2004} & виміряно \\
6 & Бістабільний нейрон: короткий вхід умикає тривале плато, гальмування вимикає; властивість умикається серотоніном (підрозділ~\ref{sec:bistable}) & Внутрішньоклітинний запис \nt{ACET}-нейронів, що керують м'язами, у кішки: коротке синаптичне збудження переводить нейрон у тривалу роботу, коротке гальмування скидає; у спінальної кішки стан з'являється після введення попередника серотоніну & \cite{Hounsgaard1988} & виміряно \\
7 & Два класи: інтегратори (частота зростає від нуля) і резонатори (на порозі одразу скінченна частота) & Класи виведено з теорії біфуркацій. Дослід: у зрізах кори \nt{GLUT}-нейрони починають працювати з як завгодно малою частотою (тип~1), швидкі \nt{GABA}-нейрони --- стрибком зі скінченної частоти (тип~2) & \cite{Izhikevich2007,Tateno2004} & виміряно \\
8 & Один нейрон --- інтегратор або детектор збігів: гальмування вкорочує вікно додавання & У зрізах гіпокампа вікно, у якому збуджувальні входи додаються до порога, коротше за $2$~мс; його вкорочує пряме гальмування, що надходить одразу за збудженням; у дендритах, де гальмування слабше, вікно ширше & \cite{Pouille2001} & виміряно \\
9 & Лінія затримки з аксонів (таблиця~\ref{tab:eetypes}) & У сипухи виміряно затримки в аксонах, що йдуть до детекторів збігів: різна довжина аксона дає різну затримку, і детектори налаштовані на різницю часів надходження звуку & \cite{Carr1990} & виміряно \\
10 & Ритмоводій під час неспання й мовчазна клітина уві сні & \nt{SERO}-нейрони працюють майже регулярно вдень, сповільнюються в повільному сні й майже мовчать у швидкому (докладніше --- розділ~\ref{sec:sleep}) & \cite{JacobsAzmitia1992,McGintyHarper1976} & виміряно \\
11 & Прилад задають іонні канали й широкомовні канали, що їх підлаштовують (твердження~\ref{prop:eetypes}) & Показано на малих мережах: речовини-регулятори змінюють власні властивості нейронів і силу синапсів, тож та сама схема видає різні ритми & \cite{Marder2012} & виміряно \\
12 & Мозок використовує переналаштування режимів для обчислень (твердження~\ref{prop:eetypes}) & Прямого досліду, де переналаштування режиму клітини було б потрібне для розв'язання задачі, немає; є лише досліди, де регулятори змінюють вихід схеми & \cite{Marder2012} & гіпотеза \\
\bottomrule
\end{longtable}}

\section{Розділ~\ref{sec:dendrites}. Кількість дендритів}

Будову класів і кількість входів виміряно анатомічно й електричними записами; оцінка~\eqref{eq:dendriteload} --- проста фізика конденсатора, а обчислювальне пояснення кількості входів спирається на теорію й моделі.

{\footnotesize
\setlength{\tabcolsep}{3pt}\renewcommand{\arraystretch}{1.15}
\begin{longtable}{>{\raggedright}p{0.5cm}>{\raggedright}p{4.0cm}>{\raggedright}p{5.6cm}>{\raggedright}p{2.3cm}>{\raggedright\arraybackslash}p{1.7cm}}
\toprule
№ & Твердження & Дослід і що виміряно & Джерело & Статус \\
\midrule\endfirsthead
\toprule
№ & Твердження & Дослід і що виміряно & Джерело & Статус \\
\midrule\endhead
1 & Питома ємність мембрани $c_m\approx1$~мкФ/см$^2$ & Із тіла нейрона знімали мембрану на піпетку, подавали сходинку напруги й за спадом ємнісного струму знаходили повну ємність, потім ділили на площу: $0.9$~мкФ/см$^2$ у трьох класів нейронів & \cite{Gentet2000} & виміряно \\
2 & Час заряджання $t\approx c_mA\Delta V/I$~\eqref{eq:dendriteload}: великому нейронові потрібні десятки одночасних входів, маленькому вистачає одного великого синапсу & Оцінка за законом заряджання конденсатора; числа ($20$ і $300$~пФ, $0.1$ і кілька нА) --- порядки величин & виведення в розділі & теорія \\
3 & Один величезний синапс дає струм у кілька наноамперів і відразу доводить маленький нейрон до порога & Одночасний запис закінчення й отримувача у зрізі: струм одного синапсу $2$--$13$~нА, кожен імпульс входу викликає надпорогову відповідь, затримка близько $0.4$~мс за $36^\circ$C; отримувач видає \nt{GLYC} & \cite{Borst1995,BorstSoria2012} & виміряно \\
4 & Нуль дендритів: у чутливих нейронів дотику й болю імпульс іде від датчика в спинний мозок повз тіло клітини & Будову (аксон, що ділиться надвоє біля тіла) встановлено анатомічно. Те, що проведення не потребує збудливості тіла, показано на перевіреній моделі такого нейрона: без збудливого тіла імпульс проходить розгалуження так само надійно & \cite{AmirDevor2003} & теорія \\
5 & Один дендрит: нюховий нейрон несе один тип рецептора й передає одну вхідну лінію & Огляд дослідів із генами рецепторів: кожен нюховий нейрон експресує один ген рецептора з великої родини & \cite{Mombaerts2004} & виміряно \\
6 & Два дендрити: у детекторів напрямку на звук входи від лівого й правого вуха сідають на різні дендрити & Анатомія й записи в ссавців, зібрані в огляді: входи від двох вух рознесено по двох протилежних дендритах & \cite{Grothe2010} & виміряно \\
7 & Рознесення входів по двох дендритах робить «І» суворішим & Показано на моделі: насичення~\eqref{eq:shunt} на одному дендриті не дає входам з одного боку дійти до порога, і детекція збігів поліпшується. Прямого досліду, де змінювали б розкладку входів, немає & \cite{AgmonSnir1998} & теорія \\
8 & Близько $7\cdot10^{10}$ маленьких \nt{GLUT}-нейронів кола навчання рухів, по $\sim4$ входи в кожного & Підрахунок клітин ізотропним фракціонуванням: у цій частині мозку людини близько $69$ млрд нейронів. Запис однієї такої клітини в живого щура: на дотик надходять пачки дискретних струмів від небагатьох входів, і клітина спрацьовує, коли вони додаються & \cite{Azevedo2009,Chadderton2004} & виміряно \\
9 & Пороговий елемент з $n$ входами розділяє не більше $P_{\max}\approx2n$ випадкових образів~\eqref{eq:cover}; для вихідного нейрона кола навчання рухів це межа $\sim2\cdot10^5$, а реалістична оцінка --- $\sim4\cdot10^4$ відповідностей & Межа --- класичний результат комбінаторної геометрії. Реалістична оцінка --- теорія ємності лише з додатними вагами й запасом на шум, застосована до виміряної кількості входів цього нейрона & \cite{Cover1965,Brunel2004} & теорія \\
10 & Змішувачі з малою кількістю входів потрібні, щоб розширити вхід; «чотири» близько до оптимуму & Теорія: вимірність подання, яку дає шар змішувачів, максимальна приблизно за тієї кількості входів, що спостерігається анатомічно; вихідна гіпотеза розширення --- у класичних працях & \cite{Marr1969,Albus1971,LitwinKumar2017} & гіпотеза \\
11 & Великі дерева: $10^4$--$10^5$ входів & Підрахунок синапсів на електронних знімках: $\approx1.7\cdot10^5$ синапсів на одному вихідному \nt{GABA}-нейроні кола навчання рухів у щура; близько $30\,000$ збуджувальних і $1\,700$ гальмівних входів на \nt{GLUT}-нейроні гіпокампа & \cite{NapperHarvey1988,Megias2001} & виміряно \\
12 & Гілки великого дерева --- окремі підсхеми з власними порогами; нейрон --- маленька багатошарова мережа & Точкова стимуляція двох місць на тонких дендритах: входи на одній гілці додаються за сигмоїдою, на різних гілках --- лінійно. У дендритах нейронів кори людини знайдено градуальні кальцієві імпульси, що дають змогу одній клітині розв'язувати лінійно нероздільну задачу. Модель: щоб відтворити один нейрон, потрібна глибока мережа & \cite{Polsky2004,Gidon2020,Beniaguev2021} & виміряно \\
13 & Дендрити-виходи: кожна гілка \nt{GABA}-нейрона сітківки --- окремий детектор напрямку & Двофотонний запис кальцію в окремих гілках: сигнал у гілці вибірковий до руху від тіла клітини до кінчика, а напруга тіла --- ні & \cite{Euler2002} & виміряно \\
14 & Кількість входів іде за задачею (твердження~\ref{prop:dendrites}) & Будову класів виміряно (рядки 3--13); що кількість входів обрано заради швидкості чи класифікації, показано лише теорією й моделями для окремих класів & \cite{LitwinKumar2017,AgmonSnir1998} & гіпотеза \\
\bottomrule
\end{longtable}}

\section{Розділ~\ref{sec:sleep}. Сон}

Режими сну, повтори й зменшення синапсів виміряно записами нейронів, мікродіалізом і електронною мікроскопією; розклад сну й повільні гойдалки описано моделями книги, а призначення сну --- здебільшого гіпотези.

{\footnotesize
\setlength{\tabcolsep}{3pt}\renewcommand{\arraystretch}{1.15}
\begin{longtable}{>{\raggedright}p{0.5cm}>{\raggedright}p{4.0cm}>{\raggedright}p{5.6cm}>{\raggedright}p{2.3cm}>{\raggedright\arraybackslash}p{1.7cm}}
\toprule
№ & Твердження & Дослід і що виміряно & Джерело & Статус \\
\midrule\endfirsthead
\toprule
№ & Твердження & Дослід і що виміряно & Джерело & Статус \\
\midrule\endhead
1 & Уві сні нейрони кори не мовчать; змінюється режим роботи & Тривалі записи нейронів кори в щурів: у повільному сні частота лишається ненульовою, періоди роботи чергуються з періодами тиші; після тривалого неспання частота вища в усіх станах, після сну --- нижча & \cite{Vyazovskiy2009} & виміряно \\
2 & \nt{ACET} високий удень і у швидкому сні, низький у повільному (таблиця~\ref{tab:sleepmodes}) & Мікродіаліз у корі вільно рухливих кішок: під час неспання викид ацетилхоліну на $100$--$175\%$ вищий, ніж у повільному сні, а у швидкому сні --- приблизно на рівні спокійного неспання & \cite{Marrosu1995} & виміряно \\
3 & \nt{NORA} і \nt{SERO} затихають у повільному сні й майже мовчать у швидкому & Запис окремих \nt{NORA}-нейронів у щурів і \nt{SERO}-нейронів у кішок за стадіями сну: частота максимальна під час неспання, нижча в повільному сні й майже нульова у швидкому & \cite{AstonJonesBloom1981,McGintyHarper1976} & виміряно \\
4 & У швидкому сні м'язи вимкнено гальмуванням через \nt{GABA} і \nt{GLYC} & У щурів вимірювали активність \nt{ACET}-нейронів жувального м'яза й подавали блокатори просто до них: параліч знімається, лише якщо заблоковано і \nt{GABA}-рецептори обох типів, і \nt{GLYC}-рецептори & \cite{BrooksPeever2012} & виміряно \\
5 & Тиск сну $S$ --- конденсатор із двома порогами~\eqref{eq:sleeppressure}, $\tau_r=18.2$~год, $\tau_d=4.2$~год & Модель двох процесів; сталі часу отримано з того, як потужність повільних хвиль ЕЕГ зростає з неспанням і падає уві сні, а форму порогів --- з часу спонтанного пробудження & \cite{Borbely1982,Daan1984} & теорія \\
6 & Машина сама виходить на $16.1$~год неспання і $8.0$~год сну (рис.~\ref{fig:twoprocess}) & Розрахунок за~\eqref{eq:sleeppressure} з гойдальними порогами, скрипт \texttt{models/sleep\_models.py}: $16.1$~год неспання і $7.9$--$8.0$~год сну & --- & модель \\
7 & Після $40$~год без сну $S=0.90$, але сон триває лише $8.8$~год: борг повертається глибиною, а не тривалістю & Модель: \texttt{models/sleep\_models.py} дає $S=0.90$ і $8.8$~год. Дослід: після $40.5$~год без сну у восьми людей у першу відновлювальну ніч потужність повільних хвиль ЕЕГ значуще вища за звичайну & \cite{Borbely1981} & модель \\
8 & Фізично $S$ --- аденозин & Мікродіаліз у кішок: позаклітинний аденозин в одній із ділянок, що керують неспанням, зростає за час неспання, зростає далі під час позбавлення сну й падає під час сну. Що $S$ --- лише аденозин, не показано & \cite{PorkkaHeiskanen1997,Dunwiddie2001} & гіпотеза \\
9 & У повільному сні кора гойдається: робота кілька сотень мс --- тиша, рідше ніж раз на секунду ($<1$~Гц, під наркозом частіше $0.3$--$0.4$~Гц) & Внутрішньоклітинний запис у кішок під наркозом: деполяризація $0.8$--$1.5$~с і тривала гіперполяризація, ритм $<1$~Гц, найчастіше $0.3$--$0.4$~Гц; те саме чергування роботи й тиші видно в природному сні (рядок~1) & \cite{Steriade1993,Vyazovskiy2009} & виміряно \\
10 & Гойдалки --- група зі зворотним зв'язком і втомою~\eqref{eq:updown}; за $b=5$ період $1.1$~с (значення моделі) і робота близько $35\%$ часу, за $b=1$ робота рівна (рис.~\ref{fig:updown}) & Розрахунок, скрипт \texttt{models/sleep\_models.py}: період $1.11$~с, частка часу роботи $0.35$ (середнє за цілим числом періодів); за $b=1$ частка роботи $1.0$. Період моделі коротший за виміряний (рядок~9) & --- & модель \\
11 & \nt{ACET} вимикає гойдалки й переводить кору в рівну роботу & Стимуляція ядра \nt{ACET}-нейронів у кішки блокувала повільний ритм у $79\%$ нейронів кори й переводила їх у рівну роботу, переважно через зникнення тривалих гіперполяризацій. Ацетилхолін пригнічує повільні калієві струми, що створюють утому & \cite{SteriadeAmzica1993,McCormick1992} & виміряно \\
12 & Спалахи ритму $7$--$15$~Гц породжує петля \nt{GLUT}--\nt{GABA} між корою й релейним вузлом & У зрізах релейного вузла зору в тхора спалахи створюються відповідними пачками релейних клітин на гальмування від сусіднього \nt{GABA}-ядра, яке вони самі збуджують & \cite{vonKrosigk1993,SteriadeMcCormickSejnowski1993} & виміряно \\
13 & Повтори дня йдуть із перемотуванням: у гіпокампі приблизно в $20$ разів швидше, у корі --- в $6$--$7$ разів & У повільному сні повторювані послідовності нейронів гіпокампа йдуть стисненими в часі. Після бігу доріжкою послідовності нейронів місця повторювалися в тому самому порядку короткими спалахами по $\sim100$~мс, стисненими приблизно в $20$ разів; у корі стиснення в $6$--$7$ разів & \cite{Nadasdy1999,LeeWilson2002,Euston2007} & виміряно \\
14 & Посилення узгодженості повторів із корою поліпшує пам'ять (причиновий зв'язок) & У щурів після навчання стимуляцією, прив'язаною до повторів, посилювали їхній зв'язок із повільними хвилями й спалахами ритму в корі: наступного дня пам'ять була високою, у контрольних --- на рівні випадку & \cite{Maingret2016} & виміряно \\
15 & Призначення повторів --- перемішане навчання повільної пам'яті & Теорія двох систем навчання й розрахунки на штучних мережах: послідовне навчання псує старе, перемішане --- ні. Прямого досліду на мозку, що повтори потрібні саме для цього, немає & \cite{McClelland1995} & гіпотеза \\
16 & Після сну синапси зменшуються пропорційно розміру, великі майже не змінюються & Тривимірна електронна мікроскопія $6920$ синапсів кори миші: площа контакту після сну на $\sim18\%$ менша, ніж після неспання, зменшення пропорційне розміру, великі стійкі синапси не змінюються & \cite{deVivo2017} & виміряно \\
17 & Головне завдання сну --- перенормування ваг & Гіпотеза синаптичного гомеостазу; рядки 1 і 16 з нею узгоджуються, але не доводять, що це головна функція & \cite{TononiCirelli2014} & гіпотеза \\
18 & Швидкий сон --- випробування моделі світу на стенді & Режим (таблиця~\ref{tab:sleepmodes}) виміряно; що мережа проганяє модель світу, --- теоретичне припущення, досліду немає & \cite{Hobson2009} & гіпотеза \\
19 & Промивання мозку уві сні: питання відкрите & Один дослід: уві сні міжклітинний простір на $60\%$ ширший і очищення швидше. Інший, з іншим способом вимірювання: уві сні й під наркозом очищення помітно повільніше. Результати суперечать один одному & \cite{Xie2013,Miao2024} & гіпотеза \\
20 & Місцевий сон: ділянки кори тварини, що не спить, ненадовго провалюються в тишу & У щурів після тривалого неспання нейрони окремої ділянки кори коротко замовкають, як у повільному сні, з повільною хвилею в місцевій ЕЕГ; у такі моменти щур частіше помиляється в задачі & \cite{Vyazovskiy2011} & виміряно \\
\bottomrule
\end{longtable}}

\section{Розділ~\ref{sec:livingmachine}. Машина з живих нейронів}

Поведінку культур на ґратках електродів виміряно в прямих дослідах із записом і стимуляцією; механізм залпів спирається на модель виснаження синапсів і досліди на мікрокультурах, а твердження про дофамін у чашці --- на поодинокі досліди, частину яких самі автори вважають лише демонстрацією.

{\footnotesize
\setlength{\tabcolsep}{3pt}\renewcommand{\arraystretch}{1.15}
\begin{longtable}{>{\raggedright}p{0.5cm}>{\raggedright}p{4.0cm}>{\raggedright}p{5.6cm}>{\raggedright}p{2.3cm}>{\raggedright\arraybackslash}p{1.7cm}}
\toprule
№ & Твердження & Дослід і що виміряно & Джерело & Статус \\
\midrule\endfirsthead
\toprule
№ & Твердження & Дослід і що виміряно & Джерело & Статус \\
\midrule\endhead
1 & Культура на чипі: здебільшого \nt{GLUT}, менша частка \nt{GABA}, що залежить від препарату; у культурах гіпокампа близько $6\%$ & Мережа з тисяч нейронів на ґратці електродів --- стандартний дослід (рядок~3). Частку \nt{GABA}-нейронів у культурах гіпокампа рахували за фарбуванням на GABA: близько $6\%$ нейронів. Для культур кори перевіреного числа ми не наводимо & \cite{Benson1994} & виміряно \\
2 & Органоїди: тривимірна нервова тканина розміром близько міліметра зі стовбурових клітин людини & Зі стовбурових клітин виростили тривимірні органоїди з кількома ділянками мозку, зокрема корою з шарами попередників і зрілими нейронами & \cite{Lancaster2013} & виміряно \\
3 & Без входу культура спалахує залпами з інтервалами від секунди до хвилин, залежно від культури & $58$ культур кори щільністю від $3\,000$ до $50\,000$ нейронів записували п'ять тижнів ($963$ півгодинні записи): після двох тижнів майже в усіх культурах переважають залпи всієї мережі; інтервали між залпами --- від $1$ до $300$~с і сильно залежать від препарату & \cite{Wagenaar2006} & виміряно \\
4 & Залп закінчується виснаженням медіатора, інтервал $T_{\text{залп}}\approx\tau_{\text{відн}}\ln\frac{1}{1-x^*}$~\eqref{eq:burstinterval}; $2$--$4$~с --- оцінка моделі з $\tau_{\text{відн}}=0.8$~с, виміряне відновлення повільніше & Формулу виведено в розділі. Модель: мережа з синапсами, що виснажуються, сама породжує залпи всієї мережі. Дослід на мікрокультурах із $4$--$30$ нейронів: тривалість залпу залежить від часу після попереднього, виснаження пухирців медіатора скасовує залпи; висновок авторів --- залп обмежений запасом медіатора. Відновлення запасу там --- десятки секунд, що з тією самою формулою дає інтервали в десятки секунд і хвилини & висновок у розділі; \cite{Tsodyks2000,CohenSegal2011,Tsodyks1997} & теорія \\
5 & «Учитель, який замовкає»: стимуляцію вимикають, коли з'явилася потрібна відповідь, і відповідь закріплюється & Культуру стимулювали з частотою $0.3$--$1$~Гц, доки через $50\pm10$~мс після стимулу не з'явиться обрана відповідь, потім стимуляцію вимикали на $5$~хв; після кількох циклів потрібна відповідь викликалася одразу; будували криві навчання & \cite{Shahaf2001} & виміряно \\
6 & Культура грає в теніс; значуще зростання серій за сеанс --- тільки за повного зворотного зв'язку & Докладно --- у розділі~\ref{sec:dishbrain} і додатку до нього: коли замість стимулу тиша, культура наприкінці сеансу теж краща, ніж без зворотного зв'язку, але з меншим ефектом; навчання від сеансу до сеансу нестійке & \cite{Kagan2022} & виміряно \\
7 & Те, що культура вчиться передбачати свій вхід, --- гіпотеза; гучні слова про «розумність» оскаржено & Принцип вільної енергії --- теоретична пропозиція; прямого досліду, що відрізняв би його від простіших пояснень, немає. Тридцять дослідників у статті-відповіді вказали, що зміна поведінки в петлі не показує ні розуму, ні відчуттів & \cite{Friston2010,Balci2023} & гіпотеза \\
8 & Культура веде віртуальне тіло до цілі за дібраного малюнка стимулів & Програма сама добирала навчальні малюнки стимуляції за поточним результатом; за десятки хвилин мережа досягала заданих станів, пластичність трималася $80$~хв після $2$~год навчання. Ті самі стимули, подані без зв'язку з поведінкою, не діяли & \cite{Bakkum2008} & виміряно \\
9 & Резервуар: навчається тільки лінійне зчитування $y=W_{\text{вих}}\mathbf r$~\eqref{eq:reservoir} & Теорія обчислення на резервуарі: будь-яка нелінійна система зі згасною пам'яттю плюс навчений лінійний вихід & \cite{Maass2002,JaegerHaas2004} & теорія \\
10 & Органоїд як резервуар розрізняє мовців із точністю близько $78\%$ (за даними авторів); за помилкою навчається зчитування, а зв'язки органоїду змінюються без учителя & Звуки мовлення кількох мовців подавали органоїду малюнками струму на ґратці електродів; навчене зчитування розрізняло мовця з точністю близько $78\%$ --- так повідомляють автори. Вони ж повідомляють, що під час тренування змінювалася функціональна зв'язність органоїду, і називають це навчанням без учителя в самому органоїді & \cite{Cai2023} & виміряно \\
11 & Мільйон нейронів витрачає на імпульси близько $10^{-4}$~Вт~\eqref{eq:dishpower} & Оцінка: ціна імпульсу $10^{-10}$~Дж з енергетичного бюджету кори~\eqref{eq:energy}, помножена на кількість імпульсів; прямого вимірювання потужності культури немає & висновок у розділі; \cite{Attwell2001} & теорія \\
12 & Дофамін за спалахом світла: $1$~мМ дофаміну в клітці, $240$ повторів, кількість викликаних імпульсів зростала & На віддаленій платформі з органоїдами дофамін у світлочутливій клітці ($1$~мМ) вивільняли ультрафіолетом ($365$~нм, $0.8$~с), коли стимул викликав більше імпульсів; за $240$ кроків (близько $5$~год) кількість імпульсів загалом зростала. Автори пишуть, що за одним дослідом зв'язок із дофаміном встановити не можна; контролів немає & \cite{Jordan2024} & гіпотеза \\
13 & Дофамін від живих клітин змінює вагу органічного транзистора надовго й оборотно & Клітини, що виділяють дофамін, вирощено над органічним нейроморфним елементом; окиснення дофаміну біля електрода змінює провідність каналу; показано довготривалу зміну ваги та її повернення & \cite{Keene2020} & виміряно \\
14 & Органоїди з робочими \nt{DOPA}-нейронами існують; як учителя їхній дофамін не використовували & В органоїдах зі стовбурових клітин людини знайдено електрично активні зрілі \nt{DOPA}-нейрони й вироблення дофаміну. Дослідів, де цей дофамін навчає іншу мережу, ми в літературі не знайшли & \cite{Jo2016} & виміряно \\
15 & Органоїд балансує жердиною: навчальні стимули від програми кращі за випадкові, без закріплення, через \nt{GLUT}-синапси & Органоїди кори миші в замкненій петлі із задачею «жердина на візку»: у більшості органоїдів сигнали, обрані навчанням із підкріпленням, дали кращий результат, ніж випадкові або жодні; після $45$~хв відпочинку покращення не зберігалося; блокада AMPA- і NMDA-рецепторів його скасовувала & \cite{Robbins2026} & виміряно \\
16 & Без регістрів керування мережа на стенді не може сама вирішити, що вважати успіхом (твердження~\ref{prop:livingmachine}) & В усіх дослідах рядків 5--15 навчальний сигнал задає експериментатор або програма; досліду, де культура сама виробила б критерій успіху, немає --- це висновок із відсутності, а не дослід & --- & гіпотеза \\
\bottomrule
\end{longtable}}

\section{Розділ~\ref{sec:dishbrain}. DishBrain}

Будову стенда й результати гри взято з однієї експериментальної роботи та її продовження; формули шуму і дрейфу до добрих налаштувань виведено в розділі й перевірено моделлю книги, а механізм навчання в чашці не встановлено.

{\footnotesize
\setlength{\tabcolsep}{3pt}\renewcommand{\arraystretch}{1.15}
\begin{longtable}{>{\raggedright}p{0.5cm}>{\raggedright}p{4.0cm}>{\raggedright}p{5.6cm}>{\raggedright}p{2.3cm}>{\raggedright\arraybackslash}p{1.7cm}}
\toprule
№ & Твердження & Дослід і що виміряно & Джерело & Статус \\
\midrule\endfirsthead
\toprule
№ & Твердження & Дослід і що виміряно & Джерело & Статус \\
\midrule\endhead
1 & Стенд: близько $800\,000$ клітин кори миші або близько $10^6$ клітин людини на чипі; $26\,000$ електродів на $8$~мм$^2$, запис до $1024$, стимуляція через $8$ & Методи роботи: чип MaxOne, $26\,000$ платинових електродів на $8$~мм$^2$, запис $1024$ каналів по $20\,000$ відліків на секунду; стимулювати технічно можна до $32$ електродів, незалежно вдалося $8$; мишачі культури використовували приблизно з $10$~діб & \cite{Kagan2022} & виміряно \\
2 & Петля з тактом $10$~мс: імпульси рухових ділянок рухають ракетку (рис.~\ref{fig:dishloop}) & Імпульси підраховують вікнами по $10$~мс; затримка від імпульсу до відповідного стимулу близько $5$~мс, її задає буфер апаратури & \cite{Kagan2022} & виміряно \\
3 & Вхід: код місця (який із $8$ електродів) і частота $4$--$40$~імп/с & В основному досліді частота стимуляції зростає лінійно від $4$~Гц, коли м'яч біля дальньої стінки, до $40$~Гц біля ракетки; амплітуда імпульсів м'яча --- $75$~мВ; імпульси двофазні прямокутні & \cite{Kagan2022} & виміряно \\
4 & Вихід $\Delta p=c\,(n_\uparrow-n_\downarrow)$~\eqref{eq:dishreadout}, шум лічби за вікно $1/\sqrt{RT}$ & Правило різниці двох ділянок описано в роботі (з вагами ділянок за активністю), точної формули не дано. Оцінка шуму --- наслідок пуассонівської лічби, виведена в розділі & \cite{Kagan2022}; висновок у розділі & теорія \\
5 & Учитель: після удару --- передбачуваний стимул $75$~мВ, $100$~імп/с, $0.1$~с; після промаху --- непередбачуваний, $150$~мВ, $5$~імп/с, $4$~с у випадкові місця й моменти, потім $4$~с паузи & Методи: після удару $75$~мВ, $100$~Гц, $100$~мс на всі $8$ електродів одразу; після промаху $150$~мВ, $5$~Гц у випадкові електроди у випадкові моменти впродовж $4$~с, потім $4$~с без стимуляції. Дофаміну в культурі немає & \cite{Kagan2022} & виміряно \\
6 & Мережа діє так, щоб її вхід став передбачуваним & Принцип вільної енергії --- теоретична пропозиція, з якої автори виводили протокол; дослід із нею узгоджується, але не відрізняє її від простіших механізмів (рядки~7--8) & \cite{Friston2010,Kagan2022} & гіпотеза \\
7 & Якщо невдача струшує мережу, а вдача ні, час у налаштуванні $\rho\propto1/P_{\text{пр}}$~\eqref{eq:dishdensity} (твердження~\ref{prop:dishscramble}) & Випливає з балансу потоків у марковському ланцюгу із симетричними поштовхами, виведено в розділі. Прямої перевірки на культурі немає & висновок у розділі & теорія \\
8 & Модель «перемішування при промаху» навчається: частка ударів $0.21\to0.38$; за поштовху після кожного розіграшу $\approx0.12$, без поштовхів $0.19$ (рис.~\ref{fig:dishmodel}) & Розрахунок, скрипт \texttt{models/dishbrain\_model.py}, $40$ моделей-культур по $2000$ розіграшів: $0.21\to0.38$; $0.13\to0.11$; $0.19\to0.19$ & --- & модель \\
9 & Серії відбитих м'ячів подовжуються тільки в живих культур у повній петлі; контролі не навчаються & $399$ сеансів по $20$~хв: $9$ культур миші ($101$ сеанс), $11$ людини ($138$), $6$ чипів із середовищем ($80$), $20$ культур у спокої ($42$), $3$ симуляції ($38$). Середня довжина розіграшу за останні $15$~хв більша, ніж за перші $5$, тільки в живих культур; у клітин, що не є нейронами (HEK293T), і в чипів із середовищем ефекту немає & \cite{Kagan2022} & виміряно \\
10 & Значуще зростання за сеанс --- тільки за повного зворотного зв'язку; коли замість стимулу тиша, наприкінці сеансу гра краща, ніж без зворотного зв'язку, але ефект менший & $486$ сеансів за $3$ дні: «стимул» ($n=27$), «тиша» ($17$), «без зворотного зв'язку» ($15$). Зростання в часі значуще тільки в умові «стимул»; наприкінці сеансу умова «тиша» перевершувала «без зворотного зв'язку» з меншим ефектом & \cite{Kagan2022} & виміряно \\
11 & Ефект невеликий; чи навчається мережа надовго --- відкрите питання & Усередині сеансу покращення надійне, але, за словами самих авторів, стійкого навчання між сеансами впродовж кількох днів не спостерігали & \cite{Kagan2022} & виміряно \\
12 & Під час гри мережа наближається до критичного режиму, і що ближче, то краща гра & Аналіз лавин активності тих самих культур: структурований вхід наближає мережу до критичного режиму, якість гри корелює з близькістю до нього; але самої лише критичності без відомостей про наслідки дій для навчання замало & \cite{Habibollahi2023} & виміряно \\
13 & «Розумність» культури не показано & Стаття-відповідь тридцяти дослідників: зміна поведінки в замкненій петлі не доводить ні розуму, ні відчуттів; досліду, який би це показав, немає & \cite{Balci2023} & гіпотеза \\
14 & Запропонований четвертий контроль: передбачуваний стимул після промаху тієї самої тривалості й енергії (підрозділ~\ref{sec:dishspec}) & Такого досліду не поставлено; це пропозиція книги & --- & гіпотеза \\
\bottomrule
\end{longtable}}

